\documentclass[10pt,journal]{IEEEtran}

\usepackage[T1]{fontenc}
\usepackage{amsmath,amssymb,amsfonts,amsthm}
\usepackage{booktabs}
\usepackage{graphicx}
\usepackage{multirow}
\usepackage{array}
\usepackage{url}
\usepackage{hyperref}
\usepackage{xcolor}
\usepackage{cite}

\newtheorem{theorem}{Theorem}
\newtheorem{proposition}{Proposition}
\newtheorem{lemma}{Lemma}
\newtheorem{definition}{Definition}

\hypersetup{
    colorlinks=true,
    linkcolor=blue!65!black,
    citecolor=blue!65!black,
    urlcolor=blue!65!black
}

\setlength{\tabcolsep}{4.2pt}
\renewcommand{\arraystretch}{1.12}

\markboth{ARXIV PREPRINT}{TIWATER RESEARCH: EVIDENCESCENE: AUTOMATED 3D SCENE RECONSTRUCTION-DRIVEN TELEOPERATION ASSISTANCE SYSTEM}

\title{EvidenceScene: An Automated 3D Scene Reconstruction-Driven Teleoperation Assistance System for Humanoid Bimanual Manipulation}

\author{Tiwater Research%
\thanks{The authors are with Tiwater Research, Tiwater (\url{https://github.com/tiwater}). Project page: \url{https://tiwater.github.io/evidencescene/}.}}

\begin{document}
\maketitle

\begin{abstract}
Remote human-in-the-loop teleoperation of humanoid bimanual manipulation tasks---such as grasping a prismatic beverage carton and pouring liquid into a narrow cylindrical cup---is severely bottlenecked by multi-body visual self-occlusion, cross-camera viewpoint discontinuity, and monocular perspective depth-parallax ambiguity when operators rely on conventional 2D camera streams. While automated 3D scene reconstruction offers a principled foundation for spatial teleoperation assistance, existing monocular depth foundation models, feed-forward multi-view pointmap transformers, and unconstrained generative mesh completion networks either exhibit decimeter-scale metric drift or hallucinate unverified geometry in safety-critical workspaces. In this paper, we present \textbf{EvidenceScene} (\emph{Observability-Conditioned 3D Scene Reconstruction for Robot Manipulation}), an end-to-end \textbf{teleoperation assistance framework and system driven by automated, sensor-grounded 3D scene reconstruction} on a 61-link humanoid mobile manipulator (16 active torso, head, and dual-arm degrees of freedom plus parallel-jaw grippers) equipped with four heterogeneous RGB-D sensors (head, chest, left-wrist, and right-wrist cameras). We first conduct a systematic empirical benchmark of \textbf{31 candidate perception, depth-estimation, feature-matching, point-tracking, and 3D reconstruction frameworks} (including Prior-DA, Depth Pro, Depth Anything V2, UniDepthV2, VGGT, SLAM3R, DUSt3R, MASt3R, CUT3R, Florence-2, Grounding DINO, SAM~2.1, YOLO11-seg, BundleTrack, BundleSDF, FoundationPose, Any6D, BIT, 6DOPE-GS, KV-Tracker, Point2Pose, Online BootsTAPIR, LoFTR, EfficientLoFTR, XFeat, ALIKED+LightGlue, and RoMa), establishing why sensor-anchored multi-camera geometric fusion outperforms standalone neural depth or pointmap inference and formalizing why increasing pose availability or lowering internal factor-graph cost can degrade strict pose accuracy and accumulated 3D surfaces under degenerate depth support. Building on these findings, our \textbf{EvidenceScene} reconstruction engine synthesizes four core technical pillars to power real-time operator assistance: (1) \emph{MoCap-free bilateral kinematic-visual extrinsic closure} on $\mathrm{SE}(3)^3$ across $1{,}648$ training frames ($12{,}061$ observations) and $413$ held-out test frames ($2{,}690$ observations) evaluated without test-split scene refitting, achieving $2.52\,\mathrm{mm}$ (left arm) and $6.56\,\mathrm{mm}$ (right arm) median 3D tabletop residuals; (2) \emph{automated 3D functional workspace reconstruction and causal pose/surface governance} in the metric tabletop frame $O_{\mathrm{table}}$ ($251{,}182$ depth inliers across $0.92\times 0.49\,\mathrm{m}$), combining watertight parametric revolution completion of the target receptacle ($\varnothing_{\mathrm{top}} = 85.4\,\mathrm{mm}$, cross-camera rim agreement $\le 1.60\,\mathrm{mm}$) with causal multi-view surfel accumulation ($5{,}699$ surfels), deterministic CUDA RANSAC and exact $\mathrm{SE}(3)$ inverse-pose transfer (raising public Laval Clock/Jar strict successes from $135/107$ to $191/136$), guarded \texttt{NO\_BA} non-keyframe retirement (saving $928.4$--$1{,}994.0\,\mathrm{MiB}$ RSS), array-native \texttt{support10/coarse12mm} $1{,}500$-point budget selection (raising public $C_{10}$ from $40.57\%/52.34\%$ to $85.10\%/85.05\%$ and task surface coverage from $71.78\%$ to $89.97\%$ while cutting selector latency by $33\%$), and initial upright bounding-prism completion ($67.8\times 94.8\times 199.5\,\mathrm{mm}$) across $498/498$ composite frames; (3) \emph{projective $\mathbb{RP}^2$ unoriented aperture geometry} (absorbing all $11$ observed $180^\circ$ SVD normal flips across $1{,}992$ camera frames with $0.000\,\mathrm{mm}$ task-metric drift) paired with \emph{role-isolated SAM~2 video segmentation memory, ViT feature caching, same-thread SAM~2 GPU-window CPU preparation overlap, and request-scoped $O(\text{runs})$ integer RLE evidence deduplication} (reducing four-camera full-chain $p_{95}$ latency by $21.7\%$ from $549.15\,\mathrm{ms}$ to $429.94\,\mathrm{ms}$ and cutting warm configured pouring HTTP endpoint latency by $27.6\%$ from $176.40\,\mathrm{ms}$ [$155.33/219.06\,\mathrm{ms}$ $p_{50}/p_{95}$] to $127.80/153.40\,\mathrm{ms}$, cutting exclusive outer object-world wrapper overhead by $72.9\%$ from $23.17$ to $6.28\,\mathrm{ms}$, and eliminating $140.58\,\mathrm{ms}$ mid-request Generation-2 GC stalls with $100\%$ bit-exact parity across $499$ online equivalence checks); and (4) an \emph{orthographic spatial teleoperation assistance interface} providing instantaneous ($0\,\mathrm{ms}$) canonical top-down ($XY$ planar bullseye) and side-elevation ($XZ$ vertical clearance) projections alongside a real-time vector bullseye HUD. Across the $37.47\,\mathrm{s}$ ($498$-frame) multi-camera manipulation study, \textbf{EvidenceScene} verifies $11.05\,\mathrm{mm}$ physical grasp-closure accuracy and guides continuous pouring alignment across a $124$-frame ($9.5\,\mathrm{s}$) bullseye window, reaching a peak radial safety margin of $+32.8\,\mathrm{mm}$ inside the $42.7\,\mathrm{mm}$ cup rim ($9.8\,\mathrm{mm}$ radial offset at $63.2^\circ$ tilt).
\end{abstract}

\begin{IEEEkeywords}
EvidenceScene, teleoperation assistance system, automated 3D scene reconstruction, humanoid bimanual manipulation, hand-eye kinematic calibration, foundation model benchmark, projective $\mathbb{RP}^2$ aperture geometry, orthographic spatial interface.
\end{IEEEkeywords}

\section{Introduction}
\label{sec:intro}

\IEEEPARstart{R}{emote} human-in-the-loop teleoperation remains the foundational paradigm for deploying humanoid and bimanual mobile manipulators in unstructured physical environments---both for direct supervisory execution of high-consequence tasks and for collecting high-precision demonstrations to train visuomotor imitation policies~\cite{chen2007human,rea2022still,billard2019trends,chi2023diffusion,zhao2023learning,sheridan1992telerobotics,fu2024mobile,qin2023anyteleop,pi2026pi07}. Despite rapid progress in robot actuation and whole-body controllers, the operational success rate of remote bimanual manipulation is fundamentally bottlenecked by the \emph{visual perception gap} between the remote physical workspace and the human operator~\cite{chen2007human,rea2022still,milgram1993argos,rakita2018autonomous,rosenberg1993virtual,bowyer2013active,su2021mixed,schwarz2017nimbro}.

The overarching objective of this work is to build an \textbf{automated 3D scene reconstruction-driven teleoperation assistance system}: a framework that continuously and automatically reconstructs the metric 3D state of the manipulation workspace---including the physical support tabletop, stationary target containers, dynamically manipulated objects, and the full 61-link robot embodiment---directly from onboard multi-camera RGB-D streams without external motion capture (MoCap) or pre-loaded object CAD models, and transforms that live 3D reconstruction into quantitative, parallax-free spatial guidance for the teleoperator (Fig.~\ref{fig:hero}).

\begin{figure*}[t]
    \centering
    \begin{minipage}[b]{0.325\textwidth}
        \centering
        \includegraphics[width=\textwidth]{../figures/07-teleop-hud-perspective-frame275.png}
        \vspace{-2pt}
        \centerline{\footnotesize (a) Perspective 3D Reconstructed Workspace}
    \end{minipage}\hfill
    \begin{minipage}[b]{0.325\textwidth}
        \centering
        \includegraphics[width=\textwidth]{../figures/08-teleop-hud-top-xy-bullseye-frame275.png}
        \vspace{-2pt}
        \centerline{\footnotesize (b) Orthographic Top-$XY$ Bullseye View}
    \end{minipage}\hfill
    \begin{minipage}[b]{0.325\textwidth}
        \centering
        \includegraphics[width=\textwidth]{../figures/09-teleop-hud-side-xz-clearance-frame275.png}
        \vspace{-2pt}
        \centerline{\footnotesize (c) Orthographic Side-$XZ$ Clearance View}
    \end{minipage}
    \caption{\textbf{Automated 3D scene reconstruction-driven teleoperation assistance at peak pouring alignment (Frame 275, $t=20.67\,\mathrm{s}$).} (a) Perspective 3D workspace integrating the 61-link humanoid manipulator, metric tabletop plane ($251{,}182$ depth inliers) with $10\,\mathrm{cm}$ major / $5\,\mathrm{cm}$ minor spatial grid and four $30\,\mathrm{mm}$ AprilTag anchors, watertight tapered target cup completion ($\varnothing_{\mathrm{top}}=85.4\,\mathrm{mm}$), and completed held beverage carton ($67.8\times 94.8\times 199.5\,\mathrm{mm}$) with live telemetry HUD and $48\times 48\,\mathrm{px}$ bullseye reticle. (b) Instantaneous Top-$XY$ orthographic projection eliminating perspective parallax along the vertical axis to verify that the pouring spout lies $9.8\,\mathrm{mm}$ from the cup center ($+32.8\,\mathrm{mm}$ inside the $42.7\,\mathrm{mm}$ rim radius). (c) Instantaneous Side-$XZ$ orthographic projection verifying $+40.5\,\mathrm{mm}$ vertical rim clearance at a $63.2^\circ$ pouring tilt.}
    \label{fig:hero}
\end{figure*}

Consider a representative bimanual liquid-pouring task executed on our 61-link humanoid manipulator: the operator must guide the right arm to grasp an upright beverage carton resting on the tabletop, lift and transport it toward a cylindrical receiving cup, and progressively tilt the carton up to $80^\circ$ so that the pouring spout $\mathbf{p}_{\mathrm{spout}}(t)$ remains strictly inside the vertical cylinder above the target cup aperture centered at $\mathbf{p}_{\mathrm{cup}}$ while maintaining positive vertical rim clearance $\Delta z(t) > 0$. When an operator relies on conventional 2D video streams from onboard RGB-D cameras, three compounding perceptual failure modes arise:
\begin{enumerate}
    \item \textbf{Severe Multi-Body Self-Occlusion:} As the right arm grasps and tilts the beverage carton, the distal forearm links, parallel-jaw gripper fingers, and the upper body of the carton itself occlude the target cup rim in the chest and wrist camera views.
    \item \textbf{Projective Depth-Parallax Ambiguity:} Under a perspective pinhole camera model $\pi_c(\mathbf{X}) = \mathbf{K}_c \mathbf{X} / X_z$, a pouring spout located at elevation $\Delta z$ above the cup rim but displaced by $\Delta x$ along the viewing ray projects to the exact same 2D pixel coordinates as the cup center whenever $\Delta x / \Delta z = \cot\alpha_c$, where $\alpha_c$ is the camera viewing elevation angle~\cite{milgram1993argos,hartley2004mvg}. Consequently, an operator viewing a 2D perspective monitor cannot reliably discern whether the spout is centered over the cup opening or $40$--$60\,\mathrm{mm}$ behind the rear rim.
    \item \textbf{Fragmented Single-View Depth Shells vs. Generative Hallucination:} Rendering raw single-frame RGB-D point clouds or foreground masks in a 3D viewer produces an incomplete, paper-thin front surface shell of the held object that lacks back walls, side depth, and tabletop support grounding~\cite{su2021mixed}. Conversely, applying unconstrained monocular depth networks~\cite{yang2024depthv2,wang2025priorda,bochkovskii2025depthpro,piccinelli2024unidepth,piccinelli2025unidepthv2,hu2024metric3dv2,bhat2023zoedepth,ke2024marigold}, feed-forward pointmap transformers~\cite{wang2024dust3r,leroy2024mast3r,wang2025vggt,liu2025slam3r,wang2025cut3r,wang2025spann3r,zhang2025monst3r}, or generative shape completion~\cite{park2019deepsdf,mildenhall2020nerf,kerbl20233dgs} risks inventing plausible-looking but metrically false geometry, violating the core safety requirement of teleoperation assistance: \emph{every rendered guidance primitive must reflect verified sensor or kinematic truth}.
\end{enumerate}

\begin{table}[t]
    \centering
    \caption{Mathematical Nomenclature and Coordinate Frame Definitions}
    \label{tab:nomenclature}
    \scriptsize
    \begin{tabular}{@{}lp{6.3cm}@{}}
        \toprule
        \textbf{Symbol} & \textbf{Definition and Metric Units} \\
        \midrule
        $\mathcal{C}$ & Four onboard RGB-D cameras: $\{\text{head}, \text{chest}, \text{left-wrist}, \text{right-wrist}\}$ \\
        $O_{\mathrm{table}}$ & Metric tabletop coordinate frame ($z_{\mathrm{table}}=0$ on support surface, $\mathrm{m}$) \\
        $O_{\mathrm{base}}$ & Robot kinematic base link frame of the 61-link manipulator ($\mathrm{m}$) \\
        $O_{\mathrm{obj}_0}$ & Canonical initial object frame of upright beverage carton at Frame 0 ($\mathrm{m}$) \\
        $\mathbf{q}_t \in \mathbb{R}^{16}$ & Synchronized active torso, head, and bilateral 7-DoF arm joint angles ($\mathrm{rad}$) \\
        $\mathbf{T}_{A\leftarrow B}$ & Rigid-body homogeneous transformation mapping frame $B$ to $A \in \mathrm{SE}(3)$ \\
        $\boldsymbol{\Theta}_{\mathrm{hw}}$ & Hardware camera mount transforms $\{\mathbf{X}_{\mathrm{chest}}, \mathbf{X}_{\mathrm{L}}, \mathbf{X}_{\mathrm{R}}\} \in \mathrm{SE}(3)^3$ \\
        $\mathcal{P}_{\mathrm{table}}$ & Tabletop depth inlier point set ($|p_z| \le 15\,\mathrm{mm}$, $|\mathcal{P}_{\mathrm{table}}|=251{,}182$) \\
        $\mathcal{S}_t \subset \mathbb{R}^3$ & Causal accumulated multi-view object surfel cloud ($|\mathcal{S}_{497}|=5{,}699$) \\
        $\mathbf{p}_{\mathrm{spout}}(t)$ & 3D position of held container pouring spout in $O_{\mathrm{table}}$ ($\mathrm{mm}$) \\
        $\mathbf{p}_{\mathrm{cup}}$ & 3D center of target cup top aperture in $O_{\mathrm{table}}$ ($\mathrm{mm}$) \\
        $R_{\mathrm{cup}}, H_{\mathrm{cup}}$ & Target cup top aperture radius ($42.7\,\mathrm{mm}$) and height ($163.2\,\mathrm{mm}$) \\
        $[\mathbf{n}] \in \mathbb{RP}^2$ & Projective equivalence class $\{+\mathbf{n}, -\mathbf{n}\}$ of unoriented rim plane \\
        $\Delta r_{xy}(t)$ & Planar Euclidean radial offset between spout and cup center ($\mathrm{mm}$) \\
        $m_{\mathrm{rim}}(t)$ & Signed radial safety margin inside cup rim: $R_{\mathrm{cup}} - \Delta r_{xy}(t)$ ($\mathrm{mm}$) \\
        $\Delta z(t)$ & Vertical elevation clearance of pouring spout above cup rim ($\mathrm{mm}$) \\
        $\theta(t)$ & Container pouring tilt angle relative to tabletop vertical ($\mathrm{deg}$) \\
        $d_{\mathrm{grip}}(t)$ & Euclidean distance from right gripper pad midpoint to carton ($\mathrm{mm}$) \\
        \bottomrule
    \end{tabular}
\end{table}

\subsection{Architectural Principles and Core Contributions}
To resolve these perceptual bottlenecks without compromising metric reliability, our architecture enforces three strict invariants across every perception, calibration, and rendering stage: (i) \emph{conservative geometric verification}, whereby every 3D surface primitive is parameterized directly by multi-camera RGB-D depth inliers, causal multi-view surfels, consensus circular rim fits, or calibrated 61-link forward kinematics; (ii) \emph{typed coordinate domain separation} across camera-local ($O_c$), fiducial ($O_{\mathrm{tag}}$), tabletop ($O_{\mathrm{table}}$), object-centric ($O_{\mathrm{obj}_0}$), and kinematic base ($O_{\mathrm{base}}$) frames; and (iii) \emph{decoupled assistance authority}, ensuring that visual guidance overlays inform the human operator without injecting unverified synthetic state into low-level joint controllers.

This paper makes six cohesive contributions (Table~\ref{tab:nomenclature}):
\begin{enumerate}
    \item \textbf{Comprehensive Empirical Benchmark of 31 Perception and Reconstruction Frameworks and Degenerate-Observation Observability Audit (Section~\ref{sec:benchmark}):} We systematically evaluate 31 contemporary depth foundation models, multi-view pointmap transformers, open-vocabulary segmenters, 6-DoF object trackers, and dense/semi-dense/sparse matchers (adding EfficientLoFTR, XFeat, and ALIKED+LightGlue) on synchronized humanoid manipulation recordings and public benchmarks (YCBInEOAT, Laval Vicon motion capture), and formalize why unanchored $\mathrm{SE}(3)$ gauge freedom, Schur-complement local rank deficiency ($r_i < 6$), and unreconciled TSDF pose versions cause higher pose availability to degrade accumulated 3D surfaces.
    \item \textbf{MoCap-Free Four-Camera to 61-Link Bilateral Kinematic Extrinsic Closure (Section~\ref{sec:closure}):} We formulate a coupled $\mathrm{SE}(3)^3$ bundle optimization over the chest and bilateral wrist camera hardware mounts through the 61-link kinematic tree, evaluated under a strict zero-refit held-out protocol across $1{,}648$ training frames ($12{,}061$ observations) and $413$ held-out test frames ($2{,}690$ observations), achieving $2.52\,\mathrm{mm}$ (left arm) and $6.56\,\mathrm{mm}$ (right arm) median 3D tabletop center residuals.
    \item \textbf{Automated 3D Functional Scene Completion, Native BundleTrack Governance, and Array-Native Budget Selection (Section~\ref{sec:completion}):} We reconstruct the manipulation workspace in $O_{\mathrm{table}}$ by fusing $251{,}182$ tabletop depth inliers ($0.92\times 0.49\,\mathrm{m}$), a watertight 49-vertex, 72-face tapered cylindrical container mesh ($\varnothing_{\mathrm{top}}=85.4\,\mathrm{mm}, \varnothing_{\mathrm{bot}}=73.0\,\mathrm{mm}, H=163.2\,\mathrm{mm}$), and causal multi-view object surfels ($5{,}699$ points) paired with a Frame-0 upright bounding prism ($67.8\times 94.8\times 199.5\,\mathrm{mm}$). In the native BundleTrack and \texttt{CausalObjectSurface} owner, we resolve CUDA RANSAC parallel-reduction non-determinism, fix \texttt{CUDACache::storeFrame} $\mathrm{SE}(3)$ inverse-coefficient truncation (raising public Clock/Jar strict successes from $135/107$ to $191/136$), enforce guarded \texttt{NO\_BA} non-keyframe retirement (reclaiming $928.4$--$1{,}994.0\,\mathrm{MiB}$ RSS), and introduce an array-native \texttt{support10/coarse12mm} $1{,}500$-point selector that raises public $C_{10}$ from $40.57\%/52.34\%$ to $85.10\%/85.05\%$ and task surface coverage from $71.78\%$ to $89.97\%$ while cutting selection latency by $33\%$.
    \item \textbf{Projective $\mathbb{RP}^2$ Unoriented Aperture Calculus (Section~\ref{sec:rp2}):} We prove mathematically and verify empirically across $1{,}992$ four-camera aperture measurements that SVD plane fits on circular container rims possess an intrinsic $\mathrm{O}(1)$ sign symmetry ($\pm\mathbf{v}_3$). Formalizing container openings on the real projective plane $\mathbb{RP}^2 = \mathbb{S}^2 / \{\pm 1\}$ absorbs all $11$ observed $180^\circ$ SVD sign flips with $0.000\,\mathrm{mm}$ task-relation drift while achieving $1.60\,\mathrm{mm}$ maximum center discrepancy and $0.53\,\mathrm{mm}$ radius discrepancy between independent head and chest cameras.
    \item \textbf{Role-Isolated Video Segmentation Memory and Bit-Exact Runtime Execution Optimization (Section~\ref{sec:runtime}):} We decouple primary and secondary SAM~2 video segmentation memory lifecycles across $2{,}322$ tracker invocations---recovering $+52$ valid secondary cup masks and eliminating $29$ background false-positive leaks---and combine observation-scoped ViT feature caching ($41/76$ cache hits) and early identity ordering ($35$ skipped refinements, reducing 4-camera full-chain $p_{95}$ by $21.7\%$ from $549.15$ to $429.94\,\mathrm{ms}$) with two bit-exact single-stream runtime adapters: \texttt{SAM2GpuWindowOverlapScope} and \texttt{RequestEvidenceDeduplicationScope} ($O(\text{runs})$ exact integer RLE arithmetic, cross-stage surface/medoid reuse, vectorized coplanar AprilTag RANSAC, and request-boundary cyclic GC control), cutting warm pouring HTTP endpoint latency by $27.6\%$ ($176.40 \to 127.80\,\mathrm{ms}$ $p_{50}$, $219.06 \to 153.40\,\mathrm{ms}$ $p_{95}$), reducing exclusive outer object-world wrapper overhead by $72.9\%$ ($23.17 \to 6.28\,\mathrm{ms}$), and eliminating $140.58\,\mathrm{ms}$ mid-request Generation-2 GC stalls with $100\%$ bit-exact parity across $499$ online checks.
    \item \textbf{Orthographic Spatial Teleoperation Interface and Bimanual Pouring Verification (Section~\ref{sec:ui}):} We prove why perspective views suffer from elevation-coupled parallax error and introduce instantaneous ($0\,\mathrm{ms}$) canonical orthographic projections (Top-$XY$ planar bullseye and Side-$XZ$ vertical clearance) paired with a real-time $48\times 48\,\mathrm{px}$ vector bullseye HUD, validating $11.05\,\mathrm{mm}$ grasp closure and a $124$-frame ($9.5\,\mathrm{s}$) bullseye pouring window peaking at $\Delta r_{xy}=9.8\,\mathrm{mm}$ ($+32.8\,\mathrm{mm}$ inside the cup rim at $63.2^\circ$ tilt).
\end{enumerate}

\section{Related Work and Framework Landscape}
\label{sec:related}

\begin{table*}[t]
    \centering
    \caption{Architectural Comparison of 3D Scene Reconstruction and Spatial Assistance Paradigms for Robot Manipulation and Teleoperation}
    \label{tab:related_comparison}
    \scriptsize
    \resizebox{\textwidth}{!}{%
    \begin{tabular}{@{}lccccccc@{}}
        \toprule
        \textbf{System / Paradigm Family} & \textbf{Multi-Camera} & \textbf{MoCap-Free} & \textbf{Watertight} & \textbf{Sensor-Grounded} & \textbf{Projective $\mathbb{RP}^2$} & \textbf{Role-Isolated} & \textbf{Orthographic} \\
        & \textbf{RGB-D Fusion} & \textbf{0-Refit FK} & \textbf{Workspace Mesh} & \textbf{(No Hallucination)} & \textbf{Aperture Math} & \textbf{Video Memory} & \textbf{XY/XZ HUD} \\
        \midrule
        2D Multi-Camera Switching~\cite{chen2007human,milgram1993argos,rakita2018autonomous} & Partial (2D only) & No & No & Yes & No & N/A & No (2D Persp.) \\
        Raw Point-Cloud VR / WebGL~\cite{rea2022still,su2021mixed,schwarz2017nimbro} & Yes (Sparse Cloud) & CAD Nominal & No (Front Shell) & Yes & No & N/A & No \\
        Monocular Depth / Pointmaps~\cite{yang2024depthv2,wang2025priorda,bochkovskii2025depthpro,piccinelli2024unidepth,piccinelli2025unidepthv2,hu2024metric3dv2,bhat2023zoedepth,ke2024marigold,wang2024dust3r,leroy2024mast3r,wang2025vggt,liu2025slam3r,wang2025cut3r,wang2025spann3r,zhang2025monst3r} & Relative / Pair & No & No ($0.11$--$79.4\,\mathrm{m}$ err.) & \textbf{No} (Scale Drift) & No & No & No \\
        Generative Shape Completion~\cite{park2019deepsdf,mildenhall2020nerf,kerbl20233dgs,liu2019softras} & Single-View & No & Yes (Hallucinated) & \textbf{No} (Unverified) & No & No & No \\
        Model-Free 6-DoF Trackers~\cite{wen2021bundletrack,wen2023bundlesdf,wen2024foundationpose,lee2025any6d,song2025bit,jin20256dopegs,elich2026kvtracker,point2pose2026} & Single-Camera & No & Object-Only & Yes & No & Single-Object & No \\
        Keypoint Affordances (kPAM)~\cite{manuelli2019kpam,qin2020keto,florence2018dense,pan2022tax} & Static Camera & Offline Calib. & Sparse Keypoints & Yes & Directed $\mathbb{S}^2$ & No & No \\
        \textbf{EvidenceScene (Ours)} & \textbf{4-Cam RGB-D} & \textbf{Yes ($\le 6.6\,\mathrm{mm}$)} & \textbf{Table+Cup+Carton} & \textbf{Strict Sensor-Grounded} & \textbf{Yes ($0.0\,\mathrm{mm}$ drift)} & \textbf{Yes ($2{,}322$ calls)} & \textbf{Yes ($0\,\mathrm{ms}$ switch)} \\
        \bottomrule
    \end{tabular}}
\end{table*}

\subsection{Teleoperation Assistance Interfaces, Virtual Fixtures, and Bimanual Manipulation}
Human factors studies in telerobotics have long established that remote operators suffer from severe depth-perception degradation, field-of-view truncation (``keyhole effect''), and high cognitive load when mentally fusing disjoint 2D camera feeds~\cite{chen2007human,rea2022still,billard2019trends,sheridan1992telerobotics}. Early supervisory displays introduced stereoscopic pointers and orthographic overlays such as ARGOS~\cite{milgram1993argos}, while virtual fixtures~\cite{rosenberg1993virtual,bowyer2013active} rendered geometric guidance rails onto operator monitors. Dynamic viewpoint selection methods~\cite{rakita2018autonomous} and mixed-reality point-cloud interfaces~\cite{su2021mixed,schwarz2017nimbro} allow free-viewpoint inspection, yet rendering perspective projections of sparse depth clouds leaves axial depth parallax unresolved along the optical viewing ray. Meanwhile, modern bimanual teleoperation systems---including ALOHA/ACT~\cite{zhao2023learning}, Mobile ALOHA~\cite{fu2024mobile}, Diffusion Policy~\cite{chi2023diffusion}, AnyTeleop~\cite{qin2023anyteleop}, and generalist Vision-Language-Action (VLA) foundation models such as $\pi_{0.7}$~\cite{pi2026pi07}---rely heavily on high-precision human demonstrations, where automated 3D scene reconstruction and orthographic spatial assistance directly reduce operator alignment error (Table~\ref{tab:related_comparison}).

\subsection{Monocular Depth Foundation Models and Multi-View 3D Pointmap Transformers}
Recent vision foundation models aim to recover dense 3D geometry directly from RGB images. Monocular relative and metric depth estimators---including Depth Anything V2~\cite{yang2024depthv2}, Prior Depth Anything (Prior-DA)~\cite{wang2025priorda}, Apple Depth Pro~\cite{bochkovskii2025depthpro}, UniDepth / UniDepthV2~\cite{piccinelli2024unidepth,piccinelli2025unidepthv2}, Metric3D~v2~\cite{hu2024metric3dv2}, ZoeDepth~\cite{bhat2023zoedepth}, and Marigold~\cite{ke2024marigold}---predict dense depth maps from single images or sparse depth priors. Concurrently, feed-forward multi-view pointmap transformers---including DUSt3R~\cite{wang2024dust3r}, MASt3R~\cite{leroy2024mast3r}, VGGT~\cite{wang2025vggt}, SLAM3R~\cite{liu2025slam3r}, CUT3R~\cite{wang2025cut3r}, Spann3R~\cite{wang2025spann3r}, and MonST3R~\cite{zhang2025monst3r}---regress 3D pointmaps directly in a shared coordinate frame. As we demonstrate in our systematic RTX~5090 benchmark (Section~\ref{sec:benchmark}), pure monocular metric depth models exhibit $0.48$--$79.4\,\mathrm{m}$ $p_{95}$ errors on close-range tabletop manipulation, and feed-forward pointmaps (VGGT, SLAM3R, DUSt3R, MASt3R) remain at $0.37$--$1.79\,\mathrm{m}$ $p_{95}$, proving that foundation depth models can only serve as hole-filling proposals (e.g., Prior-DA at $0.037\,\mathrm{m}$ chest $p_{95}$) anchored by calibrated RGB-D sensors.

\subsection{Open-Vocabulary Grounding, Promptable Video Segmentation, and Student Distillation}
Zero-shot object discovery and temporal mask tracking have been transformed by open-set detectors and promptable video segmenters, including Grounding DINO~\cite{liu2023grounding}, Florence-2~\cite{xiao2024florence2}, Segment Anything (SAM)~\cite{kirillov2023segment}, SAM~2 / SAM~2.1~\cite{ravi2024sam2}, Grounded SAM~\cite{ren2024groundedsam}, and DINO-X~\cite{ren2024dinox}, built atop Vision Transformer (ViT)~\cite{dosovitskiy2021vit}, Hiera~\cite{ryali2023hiera}, and SigLIP~\cite{zhai2023siglip} backbones. For high-rate four-camera execution, we also evaluate distilled instance segmentation students (YOLO11s-seg and YOLO11n-seg~\cite{jocher2024yolo11}) combined with staged SAM~2 encoding and Lucas-Kanade / Farneb\"ack optical flow propagation~\cite{lucas1981iterative,farneback2003two}. In our full-chain 3D reconstruction pipeline, we resolve the cross-role memory contamination and ViT tail-latency bottlenecks that arise when SAM~2 tracks multiple functional objects simultaneously (Section~\ref{sec:runtime}).

\subsection{6-DoF Object Pose Tracking, Neural/Gaussian Object Fields, and Point Tracking}
Model-based and category-level 6-DoF pose estimators such as PoseCNN~\cite{xiang2018posecnn}, DenseFusion~\cite{wang2019densefusion}, PVN3D~\cite{he2020pvn3d}, MegaPose~\cite{labbe2022megapose}, Gen6D~\cite{liu2022gen6d}, and OnePose++~\cite{he2022oneposepp} require either CAD models or multi-view reference scans. For model-free RGB-D tracking of novel objects, BundleTrack~\cite{wen2021bundletrack} optimizes a keypoint pose graph, while BundleSDF~\cite{wen2023bundlesdf} and FoundationPose~\cite{wen2024foundationpose} jointly optimize neural implicit fields or unified pose hypotheses. More recently, Any6D~\cite{lee2025any6d}, BIT (Bidirectional Iterative Tracking)~\cite{song2025bit} (incorporating SoftRas~\cite{liu2019softras}), 6DOPE-GS~\cite{jin20256dopegs} (using Gaussian Splatting~\cite{kerbl20233dgs}), KV-Tracker~\cite{elich2026kvtracker}, and Point2Pose~\cite{point2pose2026} have explored single-anchor, iterative, Gaussian, KV-cache, and point-trajectory paradigms. In parallel, long-horizon 2D point trackers---TAP-Vid~\cite{doersch2022tapvid}, TAPIR~\cite{doersch2023tapir}, Online BootsTAPIR~\cite{doersch2024bootstap}, TAPNext++~\cite{zholus2025tapnext}, and CoTracker3~\cite{karaev2024cotracker3}---alongside sparse and dense feature matchers (SIFT~\cite{lowe2004sift}, SuperGlue~\cite{sarlin2020superglue}, LightGlue~\cite{lindenberger2023lightglue}, LoFTR~\cite{sun2021loftr}, and RoMa~\cite{edstedt2024roma}) provide correspondence constraints across views.

\subsection{Dense RGB-D Mapping, Kinematic Hand-Eye Calibration, and Container Affordances}
Classical dense RGB-D reconstruction builds upon volumetric TSDF fusion (Curless and Levoy~\cite{curless1996volumetric}, KinectFusion~\cite{newcombe2011kinectfusion,izadi2011kinectfusion}, Voxel Hashing~\cite{niessner2013voxel}), surfel mapping (Point-Based Fusion~\cite{keller2013point}, ElasticFusion~\cite{whelan2015elasticfusion}, DynamicFusion~\cite{newcombe2015dynamicfusion}), and bundle-adjusted mapping (BundleFusion~\cite{dai2017bundlefusion}, BAD SLAM~\cite{schops2019badslam}, Open3D~\cite{zhou2018open3d}, ICP~\cite{besl1992icp}, FPFH~\cite{rusu2009fpfh}, RANSAC~\cite{fischler1981ransac}, COLMAP~\cite{schonberger2016sfm}, iSAM2~\cite{kaess2012isam2}, Ceres~\cite{agarwal2022ceres}, and CMA-ES~\cite{hansen2001cmaes}). For robotic container manipulation and liquid pouring~\cite{manuelli2019kpam,qin2020keto,florence2018dense,pan2022tax,pan2016robot,schenck2017visual,kennedy2019autonomous}, we prove that SVD rim plane fits exhibit $\mathrm{O}(1)$ sign indeterminacy and resolve it via projective $\mathbb{RP}^2$ directional geometry~\cite{hartley2004mvg,mardia2000directional,taubin1991estimation}. Finally, extending classical hand-eye calibration ($\mathbf{A}\mathbf{X}=\mathbf{X}\mathbf{B}$~\cite{tsai1989new,park1994robot,horaud1995hand,daniilidis1999hand,andreff1999online}), multi-camera rig factor graphs~\cite{shah2019calibration,pretto2023eomp,pachtrachai2018handeye,kalib2024,jcr2024,kcba2026}, and AprilTag fiducial PnP~\cite{olson2011apriltag,wang2016apriltag2,krogius2019apriltag3,collins2014ippe} on Lie groups~\cite{sola2018micro,barfoot2017state,huber1964robust} over ROS~2 / MCAP~\cite{macenski2022ros2,foxglove2022mcap}, we formulate a MoCap-free four-camera bilateral kinematic closure protocol with zero test-split refitting.

\section{System Architecture and Coordinate Frame Hierarchy}
\label{sec:architecture}

\begin{figure*}[t]
    \centering
    \begin{minipage}[b]{0.324\textwidth}
        \centering
        \includegraphics[width=\textwidth]{../figures/04-urdf-four-camera-measured-scene.png}
        \vspace{-2pt}
        \centerline{\scriptsize (a) Four-Camera RGB-D + 61-Link Kinematic Tree}
    \end{minipage}\hfill
    \begin{minipage}[b]{0.324\textwidth}
        \centering
        \includegraphics[width=\textwidth]{../figures/03-urdf-table-layers-full-robot-frame360.png}
        \vspace{-2pt}
        \centerline{\scriptsize (b) Whole-Body Kinematic State + Table (Frame 360)}
    \end{minipage}\hfill
    \begin{minipage}[b]{0.324\textwidth}
        \centering
        \includegraphics[width=\textwidth]{../figures/01-urdf-object-surface-frame245.png}
        \vspace{-2pt}
        \centerline{\scriptsize (c) Causal Object Surfels + Gripper Closure (Frame 245)}
    \end{minipage}
    \caption{\textbf{Multi-modal registration of the 61-link humanoid manipulator kinematic tree with four heterogeneous RGB-D cameras in the shared metric tabletop frame $O_{\mathrm{table}}$.} (a) Simultaneous projection of head, chest, left-wrist, and right-wrist RGB-D evidence alongside the 57 visual link meshes. (b) Whole-body kinematic state at Frame 360 aligned with tabletop depth inliers. (c) Close-up of causal multi-view object surface surfels ($5{,}699$ points) coupling with the right parallel-jaw gripper during pouring at Frame 245.}
    \label{fig:urdf_four_cam}
\end{figure*}

\subsection{61-Link Humanoid Manipulator and Four-Camera Sensing Topology}
Our experimental platform is a 61-link humanoid upper-body mobile manipulator whose kinematic tree comprises $61$ rigid links, $60$ joints, and $57$ visual CAD link meshes totaling $39.4\,\mathrm{MiB}$ of surface geometry (Table~\ref{tab:robot_topology}). Active articulation is governed by $16$ primary joint encoders spanning the wheeled mobile base and vertical torso column, pan-tilt head assembly, left 7-DoF manipulator arm with parallel-jaw gripper, and right 7-DoF manipulator arm with parallel-jaw gripper.

Synchronized multi-view RGB-D perception is provided by four onboard depth sensors $\mathcal{C} = \{\text{head}, \text{chest}, \text{left-wrist}, \text{right-wrist}\}$ operating at $15.0\,\mathrm{Hz}$ (median inter-frame interval $\Delta t = 66.67\,\mathrm{ms}$, joint-state synchronization skew $|\Delta t_q| = 12.1\,\mathrm{ms}$ median, $25.0\,\mathrm{ms}$ max; Fig.~\ref{fig:urdf_four_cam}):
\begin{itemize}
    \item \textbf{Head RGB-D Camera:} Mounted on the articulated pan-tilt head link, providing an elevated oblique overview of the manipulation workspace and target receptacle.
    \item \textbf{Chest RGB-D Camera:} Rigidly mounted on the upper torso structure, serving as the central workspace anchor with continuous visibility of the four tabletop AprilTag fiducials.
    \item \textbf{Left-Wrist RGB-D Camera:} Eye-in-hand sensor mounted on the distal left wrist link, observing the left manipulation hemisphere across a 7-DoF kinematic chain.
    \item \textbf{Right-Wrist RGB-D Camera:} Eye-in-hand sensor mounted on the distal right wrist link, moving dynamically with the right parallel-jaw gripper during container grasping and pouring.
\end{itemize}

\begin{table}[t]
    \centering
    \caption{Kinematic Subsystem and Onboard RGB-D Sensor Topology of the 61-Link Humanoid Manipulator}
    \label{tab:robot_topology}
    \scriptsize
    \resizebox{\columnwidth}{!}{%
    \begin{tabular}{@{}llccc@{}}
        \toprule
        \textbf{Subsystem / Sensor} & \textbf{Kinematic Chain} & \textbf{Active} & \textbf{Rigid Links /} & \textbf{Primary Role in} \\
        & \textbf{Mounting Scope} & \textbf{DoF} & \textbf{Visual Meshes} & \textbf{3D Reconstruction} \\
        \midrule
        Mobile Base \& Torso & Base $\to$ Upper Torso & $2$ & $13$ / $11$ & Global kinematic root $O_{\mathrm{base}}$ \\
        Pan-Tilt Head Chain & Torso $\to$ Head Mount & $2$ & $4$ / $4$ & Oblique workspace overview \\
        Left Manipulator Arm & Torso $\to$ Left Gripper & $7 + \text{Grip}$ & $22$ / $21$ & Left hemisphere \& loop closure \\
        Right Manipulator Arm & Torso $\to$ Right Gripper & $7 + \text{Grip}$ & $22$ / $21$ & Grasping, pouring \& closure \\
        \midrule
        \textbf{Total Embodiment} & \textbf{Full Kinematic Tree $\mathcal{U}$} & \textbf{16 + 2 Grip} & \textbf{61 / 57} & \textbf{1.40 ms live PoE FK} \\
        \midrule
        Chest RGB-D Sensor & Upper Torso Mount & Fixed & $\mathbf{X}_{\mathrm{chest}}^\star \in \mathrm{SE}(3)$ & Tabletop anchor \& object pose \\
        Head RGB-D Sensor & Pan-Tilt Head Link & $2$-DoF & Extrinsic to Torso & Dual-cam cup rim consensus \\
        Left-Wrist RGB-D & Distal Left Wrist & $7$-DoF & $\mathbf{X}_{\mathrm{L}}^\star \in \mathrm{SE}(3)$ & Bilateral $\mathrm{SE}(3)$ loop closure \\
        Right-Wrist RGB-D & Distal Right Wrist & $7$-DoF & $\mathbf{X}_{\mathrm{R}}^\star \in \mathrm{SE}(3)$ & Close-range grasp \& surfel fusion \\
        \bottomrule
    \end{tabular}}
\end{table}

\subsection{Five-Domain Coordinate Hierarchy}
At each synchronized frame index $t \in \{0, \dots, 497\}$, we enforce strict separation across five coordinate domains:
\begin{enumerate}
    \item \textbf{Sensor-Native Optical Frames ($O_c, c \in \mathcal{C}$):} Registered RGB image $I_t^{(c)} \in \mathbb{R}^{H\times W\times 3}$, metric depth map $D_t^{(c)} \in \mathbb{R}^{H\times W}$, and intrinsic matrix $\mathbf{K}_c$. Valid depth pixels back-project via $\pi_c^{-1}(u,v,z) = z\,\mathbf{K}_c^{-1}[u, v, 1]^\top \in \mathbb{R}^3$.
    \item \textbf{Local Fiducial Gauge ($O_{\mathrm{tag}}$):} Single-tag coordinate frame prior to multi-tag planar stabilization.
    \item \textbf{Measured Metric Tabletop Frame ($O_{\mathrm{table}}$):} The canonical scene reconstruction frame anchored by four coplanar $30\,\mathrm{mm}$ AprilTags and $251{,}182$ tabletop depth inliers, where $z_{\mathrm{table}} = 0$ coincides with the physical support plane and $\mathbf{n}_{\mathrm{table}} = [0, 0, -1]^\top$ points vertically upward (so elevation above the table is $h_{\mathrm{table}}(\mathbf{p}) = -p_z$).
    \item \textbf{Canonical Initial Object Frame ($O_{\mathrm{obj}_0}$):} Attached to the upright beverage carton at Frame 0, within which causal multi-view surfels $\mathcal{S}_t$ and the completed prism mesh $\mathcal{M}_{\mathrm{carton}}^{(0)}$ are rigidly defined and mapped into $O_{\mathrm{table}}$ via $\mathbf{T}_{\mathrm{table}\leftarrow \mathrm{obj}_0}(t) \in \mathrm{SE}(3)$.
    \item \textbf{Robot Base and 61-Link Kinematic Frame ($O_{\mathrm{base}}$):} Rooted at the robot base link, mapping each link $\ell \in \{1..61\}$ via forward kinematics $\mathbf{T}_{\mathrm{base}\leftarrow \ell}(\mathbf{q}_t)$ and linked into $O_{\mathrm{table}}$ via the closed camera-to-kinematic transformation $\mathbf{T}_{\mathrm{table}\leftarrow \mathrm{base}}(t)$.
\end{enumerate}

\section{MoCap-Free Four-Camera to Bilateral 61-Link Kinematic Extrinsic Closure}
\label{sec:closure}

\subsection{Product-of-Exponentials Forward Kinematics Across 61 Links}
For any link $\ell$ in the 61-link kinematic tree $\mathcal{U}$, let $\mathrm{path}(\mathrm{base}\to\ell) = ((i_1, j_1), \dots, (i_K, j_K))$ denote the ordered chain of parent-child links connected by joints $j_1, \dots, j_K$. Given the synchronized joint encoder state $\mathbf{q}_t$ at timestamp $t$, the pose of link $\ell$ in $O_{\mathrm{base}}$ is computed via recursive rigid-body composition:
\begin{equation}
    \mathbf{T}_{\mathrm{base}\leftarrow \ell}(\mathbf{q}_t) = \prod_{k=1}^{K} \mathbf{T}_{i_k \leftarrow j_k}^{\mathrm{URDF}} \, \exp\!\left( \hat{\boldsymbol{\xi}}_{j_k} \, q_{t, j_k} \right) \in \mathrm{SE}(3),
    \label{eq:poe_fk}
\end{equation}
where $\mathbf{T}_{i_k \leftarrow j_k}^{\mathrm{URDF}} \in \mathrm{SE}(3)$ is the fixed joint origin offset specified in the kinematic model and $\hat{\boldsymbol{\xi}}_{j_k} \in \mathfrak{se}(3)$ is the unit joint twist. For the parallel-jaw grippers on the left and right arms ($a \in \{L, R\}$), the 3D gripper pad midpoint in $O_{\mathrm{table}}$ is evaluated from the distal finger links $\ell_{\mathrm{grip},1}^{(a)}, \ell_{\mathrm{grip},2}^{(a)}$ with $\mathbf{o} = [0,0,0,1]^\top$:
\begin{equation}
    \mathbf{p}_{\mathrm{grip}}^{(a)}(t) = \frac{1}{2} \mathbf{T}_{\mathrm{table}\leftarrow \mathrm{base}}(t) \sum_{m=1}^{2} \mathbf{T}_{\mathrm{base}\leftarrow \ell_{\mathrm{grip},m}^{(a)}}(\mathbf{q}_t)\,\mathbf{o}.
    \label{eq:gripper_mid}
\end{equation}

\subsection{Coupled Chest-to-Wrist Loop Closure on $\mathrm{SE}(3)^3$}
Nominal CAD mounting transforms for embedded RGB-D sensors omit machining tolerances, internal optical center offsets, and bracket flexure, resulting in $30$--$80\,\mathrm{mm}$ projection discrepancies when wrist-camera depth observations are transformed through a 7-DoF arm chain. To calibrate the true hardware mounting transforms without external optical motion capture (MoCap), we formulate a coupled loop-closure optimization over the time-invariant hardware transforms:
\begin{equation}
    \boldsymbol{\Theta}_{\mathrm{hw}} = \big( \mathbf{X}_{\mathrm{chest}},\, \mathbf{X}_{\mathrm{L}},\, \mathbf{X}_{\mathrm{R}} \big) \in \mathrm{SE}(3)^3,
    \label{eq:theta_hw}
\end{equation}
where $\mathbf{X}_{\mathrm{chest}} \triangleq \mathbf{T}_{\ell_{\mathrm{torso}}\leftarrow \mathrm{chest}}$, $\mathbf{X}_{\mathrm{L}} \triangleq \mathbf{T}_{\ell_{\mathrm{L}}\leftarrow \mathrm{left\_wrist}}$, and $\mathbf{X}_{\mathrm{R}} \triangleq \mathbf{T}_{\ell_{\mathrm{R}}\leftarrow \mathrm{right\_wrist}}$.

Let $\mathcal{L} = \{\mathbf{p}_{k,m}^{\mathrm{table}} \in \mathbb{R}^3 \mid k \in \{0..3\}, m \in \{0..4\}\}$ denote the $20$ canonical 3D metric landmarks ($4$ tag centers $m=0$ and $16$ tag corners $m \in \{1..4\}$) of the four $30\,\mathrm{mm} \times 30\,\mathrm{mm}$ AprilTags affixed to the tabletop (Table~\ref{tab:apriltag_coords}). At any training frame $t \in \mathcal{D}_{\mathrm{train}}$ where the chest camera localizes the tabletop fiducial constellation via $\mathbf{T}_{\mathrm{table}\leftarrow \mathrm{chest}}(t)$ and wrist camera $w \in \{L, R\}$ observes landmark $(k,m)$ at 3D camera-local coordinate $\mathbf{y}_{t,k,m}^{(w)} \in \mathbb{R}^3$ (within pair skew $|\Delta t_{\mathrm{pair}}| \le 100.0\,\mathrm{ms}$), closing the kinematic chain through the torso link $\ell_{\mathrm{torso}}$ and distal wrist link $\ell_w$ predicts the tabletop landmark coordinate:
\begin{equation}
\begin{aligned}
    \mathbf{T}_{\mathrm{table}\leftarrow \mathrm{base}}(\mathbf{X}_{\mathrm{chest}}, \mathbf{q}_t) &= \mathbf{T}_{\mathrm{table}\leftarrow \mathrm{chest}}(t)\,\mathbf{X}_{\mathrm{chest}}^{-1} \\
    &\quad \cdot \mathbf{T}_{\mathrm{base}\leftarrow \ell_{\mathrm{torso}}}(\mathbf{q}_t)^{-1}, \\
    \hat{\mathbf{p}}_{t,k,m}^{\mathrm{table}}(\boldsymbol{\Theta}_{\mathrm{hw}}, \mathbf{q}_t) &= \mathbf{T}_{\mathrm{table}\leftarrow \mathrm{base}}(\mathbf{X}_{\mathrm{chest}}, \mathbf{q}_t) \\
    &\quad \cdot \mathbf{T}_{\mathrm{base}\leftarrow \ell_w}(\mathbf{q}_t)\,\mathbf{X}_w\,\tilde{\mathbf{y}}_{t,k,m}^{(w)}.
\end{aligned}
\label{eq:loop_closure}
\end{equation}

\begin{lemma}[$\mathfrak{se}(3)$ Left-Perturbation Jacobian of Bilateral Loop Closure]
\label{lem:jacobian}
Let $\mathbf{X}_w \mapsto \exp(\hat{\boldsymbol{\delta}}_w)\,\mathbf{X}_w$ and $\mathbf{X}_{\mathrm{chest}} \mapsto \exp(\hat{\boldsymbol{\delta}}_c)\,\mathbf{X}_{\mathrm{chest}}$ denote left-invariant Lie-algebra perturbations with $\boldsymbol{\delta}_w, \boldsymbol{\delta}_c \in \mathbb{R}^6$. Define the intermediate points $\mathbf{z}_w = \mathbf{X}_w \tilde{\mathbf{y}}_{t,k,m}^{(w)} \in \mathbb{R}^3$ and $\mathbf{z}_b = \mathbf{T}_{\ell_{\mathrm{torso}}\leftarrow \ell_w}(\mathbf{q}_t)\,\tilde{\mathbf{z}}_w \in \mathbb{R}^3$, and let $\mathbf{R}_{\mathrm{table}\leftarrow \ell_w}$ and $\mathbf{R}_{\mathrm{table}\leftarrow \mathrm{chest}}$ be the composite rotation matrices. Then the first-order differential of the residual $\mathbf{r}_{t,k,m} = \hat{\mathbf{p}}_{t,k,m}^{\mathrm{table}} - \mathbf{p}_{k,m}^{\mathrm{table}}$ satisfies:
\begin{align}
    \frac{\partial \mathbf{r}_{t,k,m}}{\partial \boldsymbol{\delta}_w} &= \mathbf{R}_{\mathrm{table}\leftarrow \ell_w} \begin{bmatrix} \mathbf{I}_3 & -[\mathbf{z}_w]_\times \end{bmatrix}, \nonumber \\
    \frac{\partial \mathbf{r}_{t,k,m}}{\partial \boldsymbol{\delta}_c} &= -\mathbf{R}_{\mathrm{table}\leftarrow \mathrm{chest}} \begin{bmatrix} \mathbf{I}_3 & -[\mathbf{X}_{\mathrm{chest}}\tilde{\mathbf{z}}_b]_\times \end{bmatrix}.
    \label{eq:jacobians}
\end{align}
\end{lemma}

Using Lemma~\ref{lem:jacobian}, we solve the Huber-robustified Levenberg-Marquardt optimization over $\mathcal{D}_{\mathrm{train}}$:
\begin{equation}
\begin{aligned}
    \boldsymbol{\Theta}_{\mathrm{hw}}^{\star} &= \arg\min_{\boldsymbol{\Theta}_{\mathrm{hw}} \in \mathrm{SE}(3)^3} \sum_{t \in \mathcal{D}_{\mathrm{train}}} \sum_{(k,m)} \rho_{\delta}\!\Big( \big\| \hat{\mathbf{p}}_{t,k,m}^{\mathrm{table}}(\boldsymbol{\Theta}_{\mathrm{hw}}, \mathbf{q}_t) \\
    &\qquad\qquad\qquad\qquad\qquad\qquad - \mathbf{p}_{k,m}^{\mathrm{table}} \big\|_2 \Big),
\end{aligned}
\label{eq:huber_opt}
\end{equation}
converging to first-order optimality $\|\nabla \mathcal{J}\|_\infty = 1.26 \times 10^{-8}$ (cost $0.0951$) on the left-arm chain and $4.66 \times 10^{-8}$ (cost $0.1195$) on the right-arm chain.

\begin{table}[t]
    \centering
    \caption{Calibrated 3D Metric Coordinates of the Four $30\,\mathrm{mm}$ Tabletop AprilTag Fiducials in $O_{\mathrm{table}}$}
    \label{tab:apriltag_coords}
    \scriptsize
    \resizebox{\columnwidth}{!}{%
    \begin{tabular}{@{}cccc@{}}
        \toprule
        \textbf{Tag ID} & \textbf{Side} & \textbf{Center $\mathbf{p}_{k,0}^{\mathrm{table}}$ (mm)} & \textbf{Planar Span $[X_{\min}, X_{\max}] \times [Y_{\min}, Y_{\max}]$ (mm)} \\
        \midrule
        Tag 0 & $29.98\,\mathrm{mm}$ & $(+4.5, -8.0, 0.0)$ & $[-11.2, +20.2] \times [-23.7, +7.7]$ \\
        Tag 1 & $29.98\,\mathrm{mm}$ & $(+13.8, -266.0, 0.0)$ & $[-2.5, +30.1] \times [-282.3, -249.7]$ \\
        Tag 2 & $29.98\,\mathrm{mm}$ & $(+572.9, -217.4, 0.0)$ & $[+556.4, +589.3] \times [-233.8, -200.9]$ \\
        Tag 3 & $29.98\,\mathrm{mm}$ & $(+555.4, +48.8, 0.0)$ & $[+539.7, +571.2] \times [+33.0, +64.5]$ \\
        \bottomrule
    \end{tabular}}
\end{table}

\subsection{Calibrated Extrinsic Transformations and Zero-Refit Held-Out Evaluation}
The converged hardware mounting transformations $\boldsymbol{\Theta}_{\mathrm{hw}}^{\star}$ (with translations in meters) recovered on the training split are:
\begin{align}
    \mathbf{X}_{\mathrm{chest}}^{\star} &= \begin{bmatrix}
        -0.0054 & -0.3226 & +0.9465 & +0.1203 \\
        -0.9999 & +0.0141 & -0.0009 & +0.0477 \\
        -0.0130 & -0.9464 & -0.3227 & +0.2062 \\
        0 & 0 & 0 & 1
    \end{bmatrix}, \label{eq:x_chest} \\
    \mathbf{X}_{\mathrm{L}}^{\star} &= \begin{bmatrix}
        +0.1215 & -0.9925 & -0.0135 & +0.0811 \\
        +0.9926 & +0.1215 & +0.0025 & -0.0577 \\
        -0.0009 & -0.0138 & +0.9999 & -0.0247 \\
        0 & 0 & 0 & 1
    \end{bmatrix}, \label{eq:x_left} \\
    \mathbf{X}_{\mathrm{R}}^{\star} &= \begin{bmatrix}
        +0.0169 & -0.9998 & +0.0144 & +0.0824 \\
        +0.9997 & +0.0171 & +0.0178 & -0.0532 \\
        -0.0180 & +0.0141 & +0.9997 & -0.0263 \\
        0 & 0 & 0 & 1
    \end{bmatrix}. \label{eq:x_right}
\end{align}

\begin{figure}[t]
    \centering
    \includegraphics[width=\columnwidth]{../figures/fig4_kinematic_closure_residuals.pdf}
    \caption{\textbf{Zero-refit held-out multi-camera kinematic closure residuals and downstream manipulation task consistency in $O_{\mathrm{table}}$.} Both the left arm ($p_{95} = 6.31\,\mathrm{mm}$ centers) and right arm ($p_{95} = 13.87\,\mathrm{mm}$ centers) pass well below the $20.0\,\mathrm{mm}$ acceptance threshold without any test-split scene parameter refitting.}
    \label{fig:closure_bars}
\end{figure}

\begin{table*}[t]
    \centering
    \caption{Train and Zero-Refit Held-Out Bilateral Kinematic-Visual Extrinsic Closure Results Across the 61-Link Manipulator}
    \label{tab:closure_full}
    \scriptsize
    \resizebox{\textwidth}{!}{%
    \begin{tabular}{@{}llcccccccc@{}}
        \toprule
        \textbf{Kinematic Chain} & \textbf{Evaluation Split} & \textbf{Frames} & \textbf{All Obs.} & \textbf{Center Obs.} & \textbf{Center $p_{50} / p_{95} / \max$} & \textbf{Corner $p_{50} / p_{95} / \max$} & \textbf{All-Point $p_{50} / p_{95} / \max$} & \textbf{Optimality} & \textbf{Acceptance} \\
        & & $N_{\mathrm{frames}}$ & $N_{\mathrm{all}}$ & $N_{\mathrm{center}}$ & (mm) [Target $p_{95} \le 20.0$] & (mm) & (mm) [Target $p_{95} \le 25.0$] & $\|\nabla\mathcal{J}\|_\infty$ & \textbf{Criterion} \\
        \midrule
        \multirow{2}{*}{\textbf{Left Arm Chain}}
        & Training Split & $839$ & $6{,}347$ & $1{,}339$ & $2.91 / 6.68 / 67.62$ & $3.58 / 8.46 / 171.39$ & $3.40 / 8.14 / 171.39$ & $1.26\times 10^{-8}$ & Converged \\
        & \textbf{Held-Out (Zero-Refit)} & $\mathbf{341}$ & $\mathbf{2{,}329}$ & $\mathbf{485}$ & $\mathbf{2.52 / 6.31 / 46.51}$ & $\mathbf{2.96 / 8.50 / 72.81}$ & $\mathbf{2.86 / 8.31 / 72.81}$ & Frozen $\boldsymbol{\Theta}_{\mathrm{hw}}^\star$ & \textbf{Passed} \\
        \midrule
        \multirow{2}{*}{\textbf{Right Arm Chain}}
        & Training Split & $809$ & $5{,}714$ & $1{,}186$ & $4.22 / 11.07 / 266.29$ & $4.24 / 12.07 / 54.28$ & $4.23 / 11.93 / 266.29$ & $4.66\times 10^{-8}$ & Converged \\
        & \textbf{Held-Out (Zero-Refit)} & $\mathbf{72}$ & $\mathbf{361}$ & $\mathbf{77}$ & $\mathbf{6.56 / 13.87 / 44.25}$ & $\mathbf{5.82 / 10.21 / 28.02}$ & $\mathbf{5.96 / 10.31 / 44.25}$ & Frozen $\boldsymbol{\Theta}_{\mathrm{hw}}^\star$ & \textbf{Passed} \\
        \bottomrule
    \end{tabular}}
\end{table*}

Table~\ref{tab:closure_full} and Fig.~\ref{fig:closure_bars} detail the training and zero-refit held-out verification results. With all hardware transforms frozen on the training split and zero scene parameters refitted on the held-out split, the left arm achieves a held-out center median error of $2.52\,\mathrm{mm}$ ($p_{95} = 6.31\,\mathrm{mm}$) across $2{,}329$ observations, and the right arm achieves $6.56\,\mathrm{mm}$ ($p_{95} = 13.87\,\mathrm{mm}$) across $361$ observations, both well within the $20.0\,\mathrm{mm}$ center and $25.0\,\mathrm{mm}$ all-point thresholds.

\section{Automated 3D Functional Scene Completion}
\label{sec:completion}

\begin{figure}[t]
    \centering
    \includegraphics[width=\columnwidth]{../figures/fig7_tabletop_fiducial_workspace_map.pdf}
    \caption{\textbf{Metric Top-Down ($XY$) workspace map of the reconstructed tabletop in $O_{\mathrm{table}}$.} Showing the stabilized $0.92\times 0.49\,\mathrm{m}$ tabletop footprint polygon, $251{,}182$ depth inlier surfels (downsampled to $650$ display surfels), the four $30\,\mathrm{mm}$ AprilTag fiducials (Tags 0--3), the initial beverage carton pose at Frame 0, the target cup aperture ($\varnothing 85.4\,\mathrm{mm}$), and the 498-frame tracked pouring spout trajectory.}
    \label{fig:workspace_map}
\end{figure}

\begin{figure*}[t]
    \centering
    \begin{minipage}[b]{0.325\textwidth}
        \centering
        \includegraphics[width=\textwidth]{../figures/01-table-cup-carton-frame001.png}
        \vspace{-2pt}
        \centerline{\footnotesize (a) Frame 1: Completed Objects + Table}
    \end{minipage}\hfill
    \begin{minipage}[b]{0.325\textwidth}
        \centering
        \includegraphics[width=\textwidth]{../figures/02-table-cup-carton-robot-frame001.png}
        \vspace{-2pt}
        \centerline{\footnotesize (b) Frame 1: Full Robot + Completed Scene}
    \end{minipage}\hfill
    \begin{minipage}[b]{0.325\textwidth}
        \centering
        \includegraphics[width=\textwidth]{../figures/03-table-cup-carton-robot-frame162.png}
        \vspace{-2pt}
        \centerline{\footnotesize (c) Frame 162: Grasp Closure ($11.05\,\mathrm{mm}$)}
    \end{minipage}
    \vspace{4pt}

    \begin{minipage}[b]{0.325\textwidth}
        \centering
        \includegraphics[width=\textwidth]{../figures/04-table-cup-carton-robot-frame245.png}
        \vspace{-2pt}
        \centerline{\footnotesize (d) Frame 245: Tilted Pouring ($36.3^\circ$)}
    \end{minipage}\hfill
    \begin{minipage}[b]{0.325\textwidth}
        \centering
        \includegraphics[width=\textwidth]{../figures/05-full-robot-scene-frame245.png}
        \vspace{-2pt}
        \centerline{\footnotesize (e) Frame 245: Whole-Body Manipulation View}
    \end{minipage}\hfill
    \begin{minipage}[b]{0.325\textwidth}
        \centering
        \includegraphics[width=\textwidth]{../figures/06-object-surface-completed-carton-frame245.png}
        \vspace{-2pt}
        \centerline{\footnotesize (f) Frame 245: Surfels + Completed Prism}
    \end{minipage}
    \caption{\textbf{Six-panel progression of automated 3D functional scene completion across the bimanual pouring task.} (a)--(b) At Frame 1 ($t=0.13\,\mathrm{s}$), the tabletop plane ($0.92\times 0.49\,\mathrm{m}$), tapered target cup ($\varnothing_{\mathrm{top}}=85.4\,\mathrm{mm}, H=163.2\,\mathrm{mm}$), upright beverage carton ($67.8\times 94.8\times 199.5\,\mathrm{mm}$), and 61-link robot model are registered in $O_{\mathrm{table}}$. (c) At Frame 162 ($t=12.20\,\mathrm{s}$), the right parallel-jaw gripper closes onto the carton with an $11.05\,\mathrm{mm}$ pad-midpoint-to-centroid residual. (d)--(f) At Frame 245 ($t=18.33\,\mathrm{s}$), causal multi-view surfels ($5{,}699$ points) and the completed carton prism tilt at $36.3^\circ$ inside the pouring bullseye cylinder above the target cup opening.}
    \label{fig:six_panel_progression}
\end{figure*}

\subsection{Tabletop Plane Stabilization and Fiducial Anchoring}
To reconstruct the physical manipulation support surface without synthetic plane guessing, valid depth pixels from the calibrated static cameras are back-projected into $O_{\mathrm{table}}$ and filtered by the support slab predicate $|p_z| \le 15.0\,\mathrm{mm}$. Across the manipulation sequence, $251{,}182$ depth inliers satisfy this support constraint within the bounded tabletop footprint:
\begin{equation}
    \mathcal{F}_{\mathrm{table}} = [-0.280, +0.644]\,\mathrm{m} \times [-0.337, +0.156]\,\mathrm{m} \times \{0\},
    \label{eq:table_footprint}
\end{equation}
spanning $0.9238\,\mathrm{m}$ in forward reach ($X_{\mathrm{table}}$) by $0.4926\,\mathrm{m}$ in lateral width ($Y_{\mathrm{table}}$), as shown in Fig.~\ref{fig:workspace_map}.

\subsection{Watertight Parametric Revolution Completion for Target Receptacles}
Dual-camera (head and chest) open-vocabulary grounding and rim depth extraction localize the target cup's circular aperture at $\mathbf{p}_{\mathrm{cup}} = [0.4071, -0.1271, -0.1632]^\top\,\mathrm{m}$ in $O_{\mathrm{table}}$, corresponding to an opening elevation of $H_{\mathrm{cup}} = 163.2\,\mathrm{mm}$ above the tabletop plane $z_{\mathrm{table}}=0$ and outer top radius $R_{\mathrm{top}} = 42.8\,\mathrm{mm}$ ($\varnothing_{\mathrm{top}} = 85.6\,\mathrm{mm}$, dual-camera consensus mean $85.4\,\mathrm{mm}$). Because overhead and oblique cameras only observe the upper rim ring ($36$ boundary points) and front exterior wall ($750$ accumulated surfels), raw depth rendering leaves the cup base ungrounded.

We construct a watertight 49-vertex, 72-face triangular surface of revolution $\mathcal{M}_{\mathrm{cup}}$ parameterized strictly by $(\mathbf{p}_{\mathrm{cup}}, R_{\mathrm{top}}, R_{\mathrm{bot}}, z_{\mathrm{table}}=0)$ with bottom radius $R_{\mathrm{bot}} = 36.6\,\mathrm{mm}$ ($\varnothing_{\mathrm{bot}} = 73.2\,\mathrm{mm}$) for $\phi_k = 2\pi k / 24$ ($k \in \{0..23\}$):
\begin{equation}
\begin{aligned}
    \mathbf{v}_k^{\mathrm{top}} &= \mathbf{p}_{\mathrm{cup}} + R_{\mathrm{top}}\big[\cos\phi_k,\, \sin\phi_k,\, 0\big]^\top, \\
    \mathbf{v}_k^{\mathrm{bot}} &= \big[p_{\mathrm{cup},x} + R_{\mathrm{bot}}\cos\phi_k,\, p_{\mathrm{cup},y} + R_{\mathrm{bot}}\sin\phi_k,\, 0\big]^\top,
\end{aligned}
\label{eq:cup_revolution}
\end{equation}
closed at the tabletop support plane by the central base vertex $\mathbf{v}_{48} = [p_{\mathrm{cup},x}, p_{\mathrm{cup},y}, 0]^\top$.

\subsection{Causal Multi-View Surfel Accumulation and Rigid Prism Completion}
For the dynamically manipulated beverage carton, let $\mathcal{S}_t \subset O_{\mathrm{obj}_0}$ denote the causal multi-view surfel cloud accumulated from Frame 0 up to Frame $t$, reaching $|\mathcal{S}_{497}| = 5{,}699$ surfels ($800$ representative display surfels exported for real-time WebGL rendering). To eliminate the ``hollow shell'' artifact under heavy right-hand occlusion, we extract the robust $5\text{th}$--$95\text{th}$ percentile bounding extents of the upright object at Frame 0 grounded to $z_{\mathrm{table}}=0$, yielding a watertight 8-vertex, 12-face, 12-edge rectangular prism $\mathcal{M}_{\mathrm{carton}}^{(0)}$ with dimensions $W \times D \times H = 67.8 \times 94.8 \times 199.5\,\mathrm{mm}$ ($[-33.7, +34.0] \times [-41.4, +53.4] \times [-62.3, +137.2]\,\mathrm{mm}$ in $O_{\mathrm{obj}_0}$) and canonical pouring spout $\mathbf{p}_{\mathrm{spout}}^{(0)} = [+15.4, +14.3, -52.5]^\top\,\mathrm{mm}$ (Fig.~\ref{fig:six_panel_progression}, Table~\ref{tab:completion_metrics}). At any frame $t$, both $\mathcal{S}_t$ and $\mathcal{M}_{\mathrm{carton}}^{(0)}$ transform rigidly into $O_{\mathrm{table}}$ via $\mathbf{T}_{\mathrm{table}\leftarrow \mathrm{obj}_0}(t) \in \mathrm{SE}(3)$.

\begin{table}[t]
    \centering
    \caption{Reconstructed 3D Functional Scene Primitives and Geometric Specifications in $O_{\mathrm{table}}$}
    \label{tab:completion_metrics}
    \scriptsize
    \resizebox{\columnwidth}{!}{%
    \begin{tabular}{@{}llcc@{}}
        \toprule
        \textbf{Scene Entity} & \textbf{Geometric Representation} & \textbf{Metric Extent} & \textbf{Surface Topology} \\
        \midrule
        Tabletop Plane & Stabilized Support Polygon & $923.8 \times 492.6\,\mathrm{mm}$ & $4$ verts, $2$ faces + grid \\
        Tabletop Inliers & Depth Slab $|z_{\mathrm{table}}| \le 15\,\mathrm{mm}$ & $251{,}182$ inliers & $650$ display surfels \\
        Fiducial Anchors & Coplanar AprilTags 0--3 & $4 \times (29.98\,\mathrm{mm})^2$ & $4$ centers, $16$ corners \\
        Target Cup & Tapered Revolution $\mathcal{M}_{\mathrm{cup}}$ & $\varnothing 85.4 / 73.0 \times 163.2\,\mathrm{mm}$ & $49$ verts, $72$ faces \\
        Beverage Carton & Upright Prism $\mathcal{M}_{\mathrm{carton}}^{(0)}$ & $67.8 \times 94.8 \times 199.5\,\mathrm{mm}$ & $8$ verts, $12$ faces \\
        Carton Surfels & Causal Surfel Cloud $\mathcal{S}_t$ & $5{,}699$ surfels & $498/498$ frames ($100\%$) \\
        Wire Snapshot & Array-Native \texttt{support10/coarse12mm} & $1{,}500$-pt budget & $85.1\%$ public $C_{10}$ / $90.0\%$ task \\
        \bottomrule
    \end{tabular}}
\end{table}

\subsection{Native BundleTrack Determinism, Guarded \texttt{NO\_BA} Retirement, and Array-Native $1{,}500$-Point Surface Selection}
\label{sec:native_bt_governance}
While the 498-frame composite workspace study (\texttt{20260804\_170929}) fuses multi-camera views offline, live and replay-portal execution on single-camera wrist/chest streams (\texttt{20260625\_173830}) requires strict fail-closed governance over our C++/CUDA BundleTrack pose owner and \texttt{CausalObjectSurface} (\texttt{RevisableSurfelMap}) surface owner:
\begin{enumerate}
    \item \textbf{Adapter Status Integrity and Deterministic CUDA RANSAC:} A two-line adapter guard prevents retained native \texttt{NO\_BA} states from being misclassified as historical candidates on cache misses (eliminating $6$ false historical promotions while preserving all $34/38$ valid task candidates). Inside \texttt{BundleTrack/src/cuda/}, parallel atomic tie-breaking and reversed count-comparison inequalities previously caused identical $13$-correspondence PCL inputs to split across $6$-inlier hypotheses differing by $18.00\,\mathrm{mm}$ and $6.22^\circ$; our \texttt{stable-winner} and count-direction repair collapses $100$ repeated native RANSAC+Procrustes calls onto a single deterministic winner.
    \item \textbf{Exact $\mathrm{SE}(3)$ Inverse-Pose Coefficient Transfer in \texttt{storeFrame}:} Replacing lossy bulk cast transfer with element-wise Eigen-to-\texttt{float4x4} inverse-pose coefficient copying in \texttt{CUDACache::storeFrame} ensures pinhole backprojection accuracy $\le 1.30\times 10^{-7}\,\mathrm{m}$. On public Vicon-referenced Laval sequences (\texttt{composed-native-retention-public-20261101}), this composed native build raises Clock-full strict ($10\,\mathrm{mm}, 10^\circ$) successes from $135$ to $\mathbf{191/204}$ (ADD AUC $0.8813 \to \mathbf{0.9135}$, $P_{10}$ $78.79\% \to 83.14\%$) and Cookiejar strict successes from $107$ to $\mathbf{136/295}$ (ADD AUC $0.7256 \to \mathbf{0.8085}$, $P_{10}$ $55.89\% \to 64.97\%$).
    \item \textbf{Guarded Contiguous-\texttt{NO\_BA} Non-Keyframe Retirement:} Under sustained tracking loss, \texttt{Bundler} previously retained all blocked frames in GPU/host memory. Invoking \texttt{Bundler::forgetFrame} on non-keyframes beyond the $\text{window\_size}+1$ budget only when the most recent $\text{window\_size}+1$ frame IDs are contiguous and all carry status \texttt{NO\_BA} (while protecting first, current, neighbor, local, keyframe, and pose-cached frames) reproduces all neural-network I/O, pose matrices, and terminal surface arrays byte-for-byte across $18$ complete $38$-input normal, stress, and recovery replays while reducing active frames from $37 \to 6$ in stress (saving $\mathbf{1{,}993.95\,\mathrm{MiB}}$ RSS) and from $23 \to 9$ in recovery (saving $\mathbf{928.38\,\mathrm{MiB}}$ RSS).
    \item \textbf{Lineage-Audited Surface Attribution, Columnar Surfel Engine, and Array-Native Budget Selection:} Tracing $1{,}512{,}357$ surfel contributions to $1{,}960{,}512$ source pixels and continuous M\"oller--Trumbore mesh rays reveals that replacing BundleTrack poses with Vicon ground-truth poses reduces forward $p_{95}$ surface error $F_{95}$ only from $25.36 \to 21.63\,\mathrm{mm}$ (Clock) and $30.93 \to 29.11\,\mathrm{mm}$ (Jar), because $52.60\% / 67.59\%$ of $>20\,\mathrm{mm}$ tail points originate within $3\,\mathrm{px}$ of mask boundaries. Requiring $\ge 2$ non-contributing views with $\ge 10^\circ$ angular separation cuts native $F_{95}$ to $12.50 / 17.33\,\mathrm{mm}$ ($P_{10} \to 90.05\% / 73.56\%$). Furthermore, upgrading \texttt{RevisableSurfelMap} to a sorted columnar Structure-of-Arrays (SoA) engine with 1D bijective $\texttt{int64}$ voxel key-packing and an array-native \texttt{support10/coarse12mm} $1{,}500$-point selector raises public $10\,\mathrm{mm}$ model coverage $C_{10}$ from $40.57\% / 52.34\%$ to $\mathbf{85.10\% / 85.05\%}$ (cutting reverse $R_{95}$ from $53.12 / 31.88\,\mathrm{mm}$ to $13.12 / 12.47\,\mathrm{mm}$), raises $552$-frame task $10\,\mathrm{mm}$ map coverage from $71.78\%$ to $\mathbf{89.97\%}$ ($p_{95}$ distance $42.92 \to 17.57\,\mathrm{mm}$), and across the 7-workload \texttt{kernel-v1} benchmark accelerates $60\mathrm{k}$-point surfel update ($73.00 \to 1.05\,\mathrm{ms}$, $\mathbf{69.8\times}$), snapshot ($20.48 \to 0.22\,\mathrm{ms}$, $\mathbf{93.0\times}$), budget selection ($15.52 \to 1.05\,\mathrm{ms}$, $\mathbf{14.8\times}$), and 3D circular rim consensus/aperture fitting ($4.45 \to 0.81\,\mathrm{ms}$, $\mathbf{5.5\times}$; $4.35 \to 1.17\,\mathrm{ms}$, $\mathbf{3.7\times}$) for an overall $\mathbf{8.65\times}$ geometric-mean kernel speed-up with $100\%$ bit-identical output digests.
\end{enumerate}

\section{Projective $\mathbb{RP}^2$ Unoriented Aperture Geometry and Task-Space Affordance Calculus}
\label{sec:rp2}

\subsection{Projective $\mathbb{RP}^2$ Invariance of Circular Aperture Planes}
Let $\mathcal{Q}_t^{(c)} = \{\mathbf{q}_i \in \mathbb{R}^3\}_{i=1}^N$ be the 3D rim points extracted from the upper boundary of the target container depth mask in camera $c$ at frame $t$. Let $\bar{\mathbf{q}} = \frac{1}{N}\sum_{i=1}^N \mathbf{q}_i$ be the sample centroid and $\mathbf{A} = [\mathbf{q}_1 - \bar{\mathbf{q}}, \dots, \mathbf{q}_N - \bar{\mathbf{q}}]^\top \in \mathbb{R}^{N\times 3}$ be the centered coordinate matrix with Singular Value Decomposition $\mathbf{A} = \mathbf{U}\boldsymbol{\Sigma}\mathbf{V}^\top$, where $\sigma_1 \ge \sigma_2 \ge \sigma_3 \ge 0$.

\begin{theorem}[Projective $\mathbb{RP}^2$ Gauge Invariance of Unoriented Aperture Planes]
\label{thm:rp2}
Assume $\sigma_2 > \sigma_3$, so the least-squares plane normal minimizing $\mathcal{E}(\mathbf{n}) = \|\mathbf{A}\mathbf{n}\|_2^2$ over $\mathbf{n} \in \mathbb{S}^2$ is one-dimensional. Then:
\begin{enumerate}
    \item The minimizer set on $\mathbb{S}^2$ is the $\mathrm{O}(1)$-orbit $\{\mathbf{v}_3, -\mathbf{v}_3\}$, where $\mathbf{v}_3 = \mathbf{V}_{:,3}$. Standard numerical SVD routines exhibit discontinuous sign inversions $\mathbf{v}_3 \mapsto -\mathbf{v}_3$ under arbitrarily small perturbations of $\mathbf{A}$ without a consistent exterior orientation oracle.
    \item Passing to the real projective quotient manifold $\mathbb{RP}^2 = \mathbb{S}^2 / \{\pm 1\}$ equipped with the geodesic distance
    \begin{equation}
        d_{\mathbb{RP}^2}([\mathbf{n}_1], [\mathbf{n}_2]) = \arccos\!\big(\min(1, |\mathbf{n}_1^\top \mathbf{n}_2|)\big) \in \left[0, \frac{\pi}{2}\right]
        \label{eq:rp2_metric}
    \end{equation}
    and orthogonal plane projector $\boldsymbol{\Pi}_{[\mathbf{n}]} = \mathbf{I}_3 - \mathbf{n}\mathbf{n}^\top$ yields a smooth, sign-invariant geometric representation.
    \item Every pouring task-space affordance metric ($\Delta r_{xy}, m_{\mathrm{rim}}, \Delta z, \theta$) and in-plane circular rim fit on $\boldsymbol{\Pi}_{[\mathbf{n}]}(\mathbf{q}_i - \bar{\mathbf{q}})$ is strictly invariant ($0.000\,\mathrm{mm}$ perturbation) under arbitrary sign flips $\mathbf{n} \mapsto -\mathbf{n}$.
\end{enumerate}
\end{theorem}
\begin{proof}
Because $\mathcal{E}(-\mathbf{n}) = \|-\mathbf{A}\mathbf{n}\|_2^2 = \|\mathbf{A}\mathbf{n}\|_2^2 = \mathcal{E}(\mathbf{n})$, both $\pm\mathbf{v}_3$ achieve the global minimum $\sigma_3^2$. Under the projection operator $\boldsymbol{\Pi}_{[\mathbf{n}]} = \mathbf{I}_3 - (-\mathbf{n})(-\mathbf{n})^\top = \mathbf{I}_3 - \mathbf{n}\mathbf{n}^\top$, the orthogonal projection of any point $\mathbf{x} \in \mathbb{R}^3$ onto the aperture plane is identical for $+\mathbf{n}$ and $-\mathbf{n}$, leaving the fitted aperture center $\mathbf{p}_{\mathrm{cup}}$, radius $R_{\mathrm{cup}}$, and tabletop-projected task offsets $(\Delta r_{xy}, m_{\mathrm{rim}}, \Delta z)$ invariant.
\end{proof}

\begin{table}[t]
    \centering
    \caption{Empirical Comparison of Directed $\mathbb{S}^2$ vs. Projective $\mathbb{RP}^2$ Aperture Plane Semantics Across $1{,}992$ Camera Measurements}
    \label{tab:rp2_audit}
    \scriptsize
    \begin{tabular}{@{}lcc@{}}
        \toprule
        \textbf{Evaluated Aperture Plane Metric} & \textbf{Directed $\mathbb{S}^2$} & \textbf{Projective $\mathbb{RP}^2$} \\
        \midrule
        Total evaluated aperture plane records & $1{,}992$ & $\mathbf{1{,}992}$ \\
        Observed SVD sign inversions ($\mathbf{n}_t^\top \mathbf{n}_{t-1} < 0$) & $11$ flips & $\mathbf{0}$ flips \\
        Maximum temporal normal jump (deg) & $180.00^\circ$ & $\mathbf{3.11^\circ}$ \\
        Dual-camera (head vs. chest) center $\max \Delta$ & $1.60\,\mathrm{mm}$ & $\mathbf{1.60\,\mathrm{mm}}$ \\
        Dual-camera (head vs. chest) radius $\max \Delta$ & $0.53\,\mathrm{mm}$ & $\mathbf{0.53\,\mathrm{mm}}$ \\
        Downstream task-metric drift under $\mathbf{n} \mapsto -\mathbf{n}$ & Discontinuous & $\mathbf{0.000\,\mathrm{mm}}$ \\
        \bottomrule
    \end{tabular}
\end{table}

\begin{figure}[t]
    \centering
    \includegraphics[width=\columnwidth]{../figures/fig6_6dof_carton_pose_and_jump_gate.pdf}
    \caption{\textbf{Tracked 6-DoF held-carton pose evolution and translation step-jump verification across all 498 frames.} (a) 3D pouring spout trajectory $(X, Y, -Z)$ in $O_{\mathrm{table}}$ converging onto the target cup $(X,Y)$ coordinates during the pouring bullseye window (Frames 226--352). (b) Roll $\phi(t)$, pitch $\theta_p(t)$, and total pouring tilt $\theta(t)$ reaching $79.4^\circ$. (c) Inter-frame 3D translation step norm $\|\Delta\mathbf{p}_t\|_2$ remaining strictly below the $120.0\,\mathrm{mm}$ step-jump bound ($\max = 24.3\,\mathrm{mm}$ during pouring, $100\%$ frame coverage).}
    \label{fig:pose_evolution}
\end{figure}

Table~\ref{tab:rp2_audit} summarizes the empirical verification of Theorem~\ref{thm:rp2} across $1{,}992$ four-camera measurements, while Fig.~\ref{fig:pose_evolution} illustrates the continuous 6-DoF object pose trajectory and bounded step-jump stability across all $498$ frames.

\section{Role-Isolated Video Segmentation Memory and Exact-Parity Runtime Optimization}
\label{sec:runtime}

\begin{figure*}[t]
    \centering
    \includegraphics[width=0.92\textwidth]{../figures/fig3_runtime_latency_breakdown.pdf}
    \caption{\textbf{Runtime stage latency profile and exact-parity optimization ablation across the 498-frame manipulation sequence.} (a) Per-stage median ($p_{50}$) and tail ($p_{95}$) execution latency, demonstrating that SAM~2 primary tracking ($100.22\,\mathrm{ms}$ $p_{50}$, $322.58\,\mathrm{ms}$ $p_{95}$) and pose/surfel accumulation ($75.94\,\mathrm{ms}$ $p_{50}$, $100.69\,\mathrm{ms}$ $p_{95}$) dominate compute while 61-link URDF forward kinematics requires only $1.40\,\mathrm{ms}$. (b) Exact-parity optimization ablation reducing full-chain $p_{95}$ service latency by $21.7\%$ (from $549.15\,\mathrm{ms}$ to $429.94\,\mathrm{ms}$) with $100.0\%$ trajectory parity.}
    \label{fig:latency_ablation}
\end{figure*}

\begin{table*}[t]
    \centering
    \caption{Per-Stage Runtime Latency Decomposition and Exact-Parity Optimization Ablation Across 498 Synchronized Multi-Camera Pairs}
    \label{tab:runtime_ablation}
    \scriptsize
    \resizebox{\textwidth}{!}{%
    \begin{tabular}{@{}lccccccc@{}}
        \toprule
        \textbf{Pipeline Stage / Optimization Variant} & \textbf{Service $p_{50}$} & \textbf{Service $p_{95}$} & \textbf{Service Max} & \textbf{Wall Time} & \textbf{ViT Encoder} & \textbf{Skipped Feature} & \textbf{Trajectory} \\
        & (ms) & (ms) & (ms) & (s) & \textbf{Cache Hits} & \textbf{Refinements} & \textbf{Parity} \\
        \midrule
        \multicolumn{8}{@{}l}{\emph{Part A: Per-Stage Execution Latency Breakdown (Normalized Full Pipeline)}} \\
        \quad Stage 1: SAM~2 Primary Mask Tracking & $100.22$ & $322.58$ & $612.40$ & $39.42$ & --- & --- & Primary Segmentation \\
        \quad Stage 2: 6-DoF Pose \& Surfel Accumulation & $75.94$ & $100.69$ & $148.20$ & $24.18$ & --- & --- & Object Pose \& Surfels \\
        \quad Stage 3: 3D Aperture Rim Extraction \& SVD Fit & $21.39$ & $32.21$ & $49.80$ & $6.84$ & --- & --- & Receptacle Rim Fit \\
        \quad Stage 4: Multi-Camera RGB-D Stream Decoding & $13.14$ & $19.81$ & $31.50$ & $3.12$ & --- & --- & Sensor Ingestion \\
        \quad Stage 5: Task Graph \& Secondary Renewal & $4.57$ & $7.62$ & $184.60$ & $1.47$ & --- & --- & Task Affordance DAG \\
        \quad Stage 6: 61-Link PoE Forward Kinematics & $1.40$ & $2.80$ & $4.90$ & $0.70$ & --- & --- & Kinematic Embodiment \\
        \midrule
        \multicolumn{8}{@{}l}{\emph{Part B: End-to-End Exact-Parity Optimization Ablation ($N=498$ Synchronized Pairs, $2{,}322$ Tracker Invocations)}} \\
        \quad 1. Baseline Full Causal Pipeline & $210.94$ & $549.15$ & $1{,}124.80$ & $78.57$ & $0 / 76$ ($0.0\%$) & $0$ & $498/498$ ($100.0\%$) \\
        \quad 2. + Observation-Scoped ViT Feature Cache & $215.72$ & $447.86$ & $1{,}018.40$ & $76.49$ & $41 / 76$ ($53.9\%$) & $0$ & $498/498$ ($100.0\%$) \\
        \quad 3. + Early Secondary Identity \& Role Isolation & $\mathbf{214.00}$ & $\mathbf{429.94}$ & $\mathbf{1{,}007.22}$ & $\mathbf{75.03}$ & $\mathbf{41 / 76}$ ($\mathbf{53.9\%}$) & $\mathbf{35}$ & $\mathbf{498/498}$ ($\mathbf{100.0\%}$) \\
        \midrule
        \multicolumn{8}{@{}l}{\emph{Part C: Bit-Exact Single-Stream HTTP Endpoint Optimization Ablation on Configured Pouring Sequence ($N=40$ Frames, $78$ Warm Requests / Condition)}} \\
        \quad 0. Per-Request Portal Baseline (\texttt{2026-11-01}) & $176.40$ & $233.66$ & $1{,}464.50$ & $8.32$ & $0$ hash threads saved & $735/65/3$ GC & $34$ cand / $32$ surf ($100\%$) \\
        \quad 1. Persistent Portal + Pre-Overlap Control & $155.33$ & $219.06$ & $1{,}422.12$ & $7.49$ & $0.00\,\mathrm{ms}$ GPU-win prep & $735/65/3$ GC & $400/400$ masks ($100\%$) \\
        \quad 2. + \texttt{SAM2GpuWindowOverlapScope} (\texttt{2026-11-02}) & $149.01$ & $212.24$ & $1{,}416.84$ & $7.25$ & $5.72\,\mathrm{ms}$ in $24.2\,\mathrm{ms}$ GPU wait & $38/38$ hash thr. $\to 0$ & $190/190$ native hashes \\
        \quad 3. + \textbf{\texttt{RequestEvidenceDeduplicationScope}} (\texttt{2026-11-02}) & $\mathbf{127.80}$ & $\mathbf{153.40}$ & $\mathbf{1{,}152.96}$ & $\mathbf{6.13}$ & \textbf{Wrapper $\mathbf{23.17 \to 6.28\,\mathrm{ms}}$ ($\mathbf{-72.9\%}$)} & \textbf{$\mathbf{0}$ Mid-Req Gen-2 GC} & $\mathbf{499/499}$ \textbf{Checks ($\mathbf{100\%}$)} \\
        \bottomrule
    \end{tabular}}
\end{table*}

\subsection{Role-Isolated SAM~2 Video Segmentation Sessions}
When the primary manipulated object (beverage carton) and secondary receptacle (receiving cup) share a single stateful SAM~2 predictor instance, periodic grounding renewals on the secondary role mutate shared memory-bank tensors. Formally, let $\mathcal{M}_t^{(r)} = \{(\mathbf{K}_\tau^{(r)}, \mathbf{V}_\tau^{(r)})\}_{\tau \in \mathcal{H}_t}$ denote the spatial memory bank for functional role $r \in \{\text{primary}, \text{secondary}\}$. In a shared memory bank, cross-attention over concatenated keys $[\mathbf{K}_t^{(\mathrm{pri})} \;\|\; \mathbf{K}_t^{(\mathrm{sec})}]$ introduces cross-role feature drift whenever one object occludes the other. We isolate the video segmentation state into independent per-role predictor sessions ($\mathcal{T}_{\mathrm{primary}} \perp \mathcal{T}_{\mathrm{secondary}}$) scoped by camera stream and functional role. Across all $2{,}322$ tracker invocations in the 498-frame sequence, role isolation preserves $100.0\%$ bit-exact primary trajectory parity ($498/498$ frames) while recovering $+52$ valid secondary cup masks and eliminating $29$ background false-positive feature updates.

\subsection{Observation-Scoped ViT Feature Caching and Early Identity Gating}
To eliminate the $p_{95} = 549.15\,\mathrm{ms}$ tail-latency spikes triggered when secondary role renewal re-evaluates an RGB-D frame already encoded by the primary tracker, we introduce two exact-parity systems optimizations (Fig.~\ref{fig:latency_ablation}, Table~\ref{tab:runtime_ablation}, Part~B):
\begin{enumerate}
    \item \textbf{Observation-Scoped ViT Feature Reuse:} Keyed by camera identifier and sensor timestamp, we cache the SAM~2 Hiera/ViT image embedding and Grounding DINO backbone features within the active observation scope, yielding $41/76$ ($53.9\%$) ViT encoder cache hits and reducing full-chain $p_{95}$ latency from $549.15\,\mathrm{ms}$ to $447.86\,\mathrm{ms}$ ($-18.4\%$).
    \item \textbf{Early Secondary Identity Ordering:} By evaluating the SigLIP visual-semantic similarity and geometric support criterion \emph{before} invoking SAM~2 mask feature refinement during secondary role renewal, $35$ rejected candidate proposals bypass high-resolution feature extraction entirely, reducing failed-attempt call latency from $548.6 / 577.0\,\mathrm{ms}$ ($p_{50}/p_{95}$) to $413.8 / 447.5\,\mathrm{ms}$ and bringing full-chain $p_{95}$ latency down to $\mathbf{429.94\,\mathrm{ms}}$ ($\mathbf{-21.7\%}$).
\end{enumerate}

\subsection{Same-Thread SAM~2 GPU-Window Overlap and $O(\text{runs})$ Request Evidence Deduplication}
\label{sec:gpu_window_and_dedup}
On the synchronous single-stream \texttt{/world\_model/observe} HTTP route (\texttt{teleop\_20260625\_173830}, Table~\ref{tab:runtime_ablation}, Part~C), fine-grained critical-path attribution showed that sequential SAM~2 and native BundleTrack execution ($95.87 / 118.98\,\mathrm{ms}$ warm $p_{50}/p_{95}$) was compounded by $\sim 48\,\mathrm{ms}$ of redundant host CPU preparation and mid-request cyclic garbage collection stalls. We introduce two stacked, bit-exact runtime execution adapters:
\begin{enumerate}
    \item \textbf{Same-Thread SAM~2 GPU-Window CPU Preparation Overlap (\texttt{SAM2GpuWindowOverlapScope}):} On warm distorted frames, \texttt{SAM2VideoMemoryTracker.track} enqueues CUDA kernels for \texttt{propagate\_in\_video} in $\sim 9.0\,\mathrm{ms}$ and then blocks the host thread for $\sim 24.17\,\mathrm{ms}$ at \texttt{(logits > 0.0).cpu().numpy()} waiting for GPU \texttt{memory\_attention}. Rather than spawning background \texttt{ThreadPoolExecutor} workers that contend for the Python GIL, our adapter performs a $0.25\,\mathrm{ms}$ fail-closed pre-validation, launches \texttt{propagate\_in\_video}, and executes mask-independent CPU preparation \emph{on the same thread inside the $24.17\,\mathrm{ms}$ GPU wait window} before calling \texttt{.cpu().numpy()}: (i) bilinear/nearest undistortion \texttt{rectify\_registered\_rgbd}, (ii) dual SHA-256 evidence-hash prefixes covering \texttt{header+color\_rgb} and \texttt{header+color\_rgb+depth\_mm} ($4.61 / 5.53\,\mathrm{MB} = 83.3\%$ of input bytes), and (iii) private \texttt{color\_bgr} and \texttt{depth\_m} buffers for \texttt{PinnedBundleTrackSession}. After GPU mask readback, it reuses in-scope \texttt{sensor\_mask} RLE/SHA-256 and finalizes the remaining $16.7\%$ hash tail synchronously in $0.374\,\mathrm{ms}$, absorbing $5.72\,\mathrm{ms}$ of CPU work into the GPU wait ($24.17 \to 18.45\,\mathrm{ms}$), eliminating all $38$ per-frame background hash threads, and lowering \texttt{registry\_advance\_ms} $p_{50}/p_{95}$ from $102.29/131.43\,\mathrm{ms}$ to $96.53/124.48\,\mathrm{ms}$.
    \item \textbf{Request-Scoped Evidence Deduplication and Exact $O(\text{runs})$ RLE Arithmetic (\texttt{RequestEvidenceDeduplicationScope}):} \texttt{FoundationMaskCacheDetector} previously iterated pixel-by-pixel in Python over $80{,}696$ foreground pixels per mask, allocating $>161{,}000$ Python \texttt{int} objects per frame ($10.38\,\mathrm{ms}$). Given RLE runs $(s_k, L_k)$ with $e_k = s_k + L_k - 1$ on an $H\times W$ grid, let $(y_0, x_0) = \operatorname{divmod}(s_k, W)$ and $(y_1, x_1) = \operatorname{divmod}(e_k, W)$. Using exact arbitrary-precision Python integers, the row sum $\sum y$ across a multi-row run ($y_0 < y_1$) is evaluated in $O(1)$ closed form:
    \begin{equation}
        \Sigma_{y,k} = (W - x_0)y_0 + W\frac{(y_1 - y_0 - 1)(y_0 + y_1)}{2} + (x_1 + 1)y_1,
        \label{eq:monograph_rle_sum_y}
    \end{equation}
    and the column sum follows from linear index sums $\Sigma_{\mathrm{idx},k} = L_k s_k + L_k(L_k - 1)/2$ via $\sum x = \sum_k \Sigma_{\mathrm{idx},k} - W\sum_k \Sigma_{y,k}$, cutting bounding-box and centroid evaluation from $10.38\,\mathrm{ms}$ to $0.015\,\mathrm{ms}$ ($99.85\%$ reduction) with bit-for-bit floating-point identity. Combining this with request-scoped memoization of \texttt{RegisteredRgbdFrame.select\_surface(mask)} and medoid depth across \texttt{table\_contexts}, \texttt{detector}, and \texttt{run\_live} ($41/41$ hits), vectorized $240$-hypothesis coplanar AprilTag RANSAC ($4.47 \to 0.89\,\mathrm{ms}$), and post-init \texttt{gc.collect()} plus request-boundary cyclic GC control (reclaiming $\sim 745\,\mathrm{MB}$ of init garbage, cutting CUDA allocated memory from $813\,\mathrm{MB}$ to $67.1\,\mathrm{MB}$, and eliminating $140.58\,\mathrm{ms}$ mid-request Generation-2 GC stalls) cuts exclusive outer \texttt{original\_world} wrapper overhead by $\mathbf{72.9\%}$ ($23.17 \to 6.28\,\mathrm{ms}$) and lowers warm HTTP \texttt{endpoint\_ms} $p_{50}/p_{95}$ to $\mathbf{127.80 / 153.40\,\mathrm{ms}}$ ($\mathbf{-27.54 / -65.66\,\mathrm{ms}}$ vs. pre-overlap control; $\mathbf{-27.6\%}$ vs. the $176.40\,\mathrm{ms}$ portal baseline) with $100\%$ byte-for-byte parity across $400$ SAM masks, $190$ native BundleTrack inputs, and $499$ online equivalence checks.
\end{enumerate}

\section{Orthographic Spatial Teleoperation Interface and Bimanual Pouring Evaluation}
\label{sec:ui}

\subsection{Perspective Parallax Ambiguity vs. Orthographic Decoupling}
\begin{proposition}[Perspective Parallax Error vs. Orthographic Invariance]
\label{prop:parallax}
Let $\mathbf{p}_{\mathrm{spout}} = \mathbf{p}_{\mathrm{cup}} + [\Delta x, \Delta y, -\Delta z]^\top$ in $O_{\mathrm{table}}$. Under a perspective camera viewing the cup center at elevation angle $\alpha_c \in (0, \pi/2)$ along the $XZ$ plane, the projected 2D apparent planar offset on the image plane is perturbed by the elevation-coupled parallax term
\begin{equation}
    \delta_{\mathrm{parallax}}(\Delta z, \alpha_c) = \Delta z \cot\alpha_c.
    \label{eq:parallax_bias}
\end{equation}
For $\Delta z = 81.7\,\mathrm{mm}$ (Frame 245) and $\alpha_c = 55^\circ$, perspective foreshortening induces a $\delta_{\mathrm{parallax}} = 57.2\,\mathrm{mm}$ apparent horizontal shift ($1.34\times$ the cup radius $R_{\mathrm{cup}}$). Conversely, under canonical orthographic projections $\boldsymbol{\Pi}_{\mathrm{TopXY}}(x,y,z) = (x,y)$ and $\boldsymbol{\Pi}_{\mathrm{SideXZ}}(x,y,z) = (x, -z)$, planar radial offset $\Delta r_{xy}$ and vertical clearance $\Delta z$ decouple into orthogonal, parallax-free coordinates ($\delta_{\mathrm{parallax}} \equiv 0$).
\end{proposition}
\begin{proof}
In the local camera frame tilted at elevation $\alpha_c$ above the tabletop $XY$ plane, a vertical displacement $-\Delta z\,\mathbf{e}_z^{\mathrm{table}}$ projects onto the image horizontal axis with component $\Delta z \cos\alpha_c$ while scaling the optical depth by $\sin\alpha_c$, producing a perspective ray offset equivalent to a tabletop shift of $\Delta x = \Delta z \cot\alpha_c$. Orthographic projection along $\mathbf{e}_z^{\mathrm{table}}$ ($\boldsymbol{\Pi}_{\mathrm{TopXY}}$) has nullspace $\mathrm{span}\{\mathbf{e}_z^{\mathrm{table}}\}$, annihilating $\Delta z$ identically ($\partial \boldsymbol{\Pi}_{\mathrm{TopXY}} / \partial z = \mathbf{0}$), while $\boldsymbol{\Pi}_{\mathrm{SideXZ}}$ projects orthogonally onto the vertical elevation plane.
\end{proof}

\begin{figure}[t]
    \centering
    \includegraphics[width=\columnwidth]{../figures/fig1_pouring_trajectory_telemetry.pdf}
    \caption{\textbf{Continuous 3D task-functional telemetry across all 498 synchronized frames ($37.47\,\mathrm{s}$).} (a) Planar radial alignment error $\Delta r_{xy}(t)$ relative to the $42.7\,\mathrm{mm}$ cup rim radius, highlighting the $124$-frame pouring bullseye window (Frames 226--352). (b) Vertical rim clearance $\Delta z(t)$ and container pouring tilt $\theta(t)$. (c) Euclidean distance $d_{\mathrm{grip}}(t)$ between the 61-link right gripper pad midpoint and the visually tracked carton centroid.}
    \label{fig:telemetry_full}
\end{figure}

\begin{figure*}[t]
    \centering
    \begin{minipage}[b]{0.485\textwidth}
        \centering
        \includegraphics[width=\textwidth]{../figures/fig2_orthographic_bullseye_locus.pdf}
        \vspace{-4pt}
        \centerline{\footnotesize (a) Orthographic Top-Down ($XY$) Spout Locus vs. Cup Rim}
    \end{minipage}\hfill
    \begin{minipage}[b]{0.485\textwidth}
        \centering
        \includegraphics[width=\textwidth]{../figures/fig5_side_xz_pouring_profile.pdf}
        \vspace{-4pt}
        \centerline{\footnotesize (b) Orthographic Side-Elevation ($XZ$) Spout Profile \& Cup Wall}
    \end{minipage}
    \caption{\textbf{Orthographic spatial decomposition of the 498-frame pouring trajectory in $O_{\mathrm{table}}$.} (a) In the Top-$XY$ plane, the operator guides the pouring spout into the $R=42.7\,\mathrm{mm}$ aperture circle at Frame 226 and maintains bullseye alignment through Frame 352, reaching $9.8\,\mathrm{mm}$ radial offset ($+32.8\,\mathrm{mm}$ safety margin) at Frame 275. (b) In the Side-$XZ$ elevation plane, the spout approaches over the completed $163.2\,\mathrm{mm}$-tall tapered cup wall, maintaining $+81.7\,\mathrm{mm}$ vertical rim clearance at entry (Frame 245), $+40.5\,\mathrm{mm}$ at peak bullseye (Frame 275, $63.2^\circ$ tilt), and $+33.3\,\mathrm{mm}$ at maximum tilt (Frame 315, $79.4^\circ$ tilt).}
    \label{fig:xy_xz_pair}
\end{figure*}

\begin{figure}[t]
    \centering
    \includegraphics[width=\columnwidth]{../figures/fig8_phase_space_margin_vs_tilt.pdf}
    \caption{\textbf{Manipulation phase-space safety diagram: signed radial rim margin $m_{\mathrm{rim}}(t) = R_{\mathrm{cup}} - \Delta r_{xy}(t)$ vs. container pouring tilt angle $\theta(t)$ and vertical clearance $\Delta z(t)$.} The trajectory transitions cleanly from standby transport ($m_{\mathrm{rim}} < 0, \theta < 25^\circ$) into the safe aperture cylinder ($m_{\mathrm{rim}} \ge 0$) prior to high-angle tilting ($\theta \ge 45^\circ$), strictly avoiding the spill-risk forbidden quadrant ($m_{\mathrm{rim}} < 0, \theta \ge 45^\circ$).}
    \label{fig:phase_space}
\end{figure}

\subsection{Quantitative Analysis of Bimanual Grasp Closure and Bullseye Pouring}
Guided by Proposition~\ref{prop:parallax}, our spatial assistance interface provides instantaneous ($0.0\,\mathrm{ms}$, $60\,\mathrm{fps}$) camera-projection switching between Isometric 3D, Top-$XY$ (Bullseye), and Side-$XZ$ (Clearance) views, accompanied by a live $48\times 48\,\mathrm{px}$ vector bullseye reticle inside the telemetry HUD (Fig.~\ref{fig:hero}, Table~\ref{tab:modality_comparison}).

Figures~\ref{fig:telemetry_full}, \ref{fig:xy_xz_pair}, and \ref{fig:phase_space} and Table~\ref{tab:phase_summary} document the quantitative performance across the five physical manipulation phases of the 498-frame ($37.47\,\mathrm{s}$) bimanual pouring task:
\begin{enumerate}
    \item \textbf{Phase I --- Pre-Grasp Reach (Frames 0--139, $0.00$--$10.53\,\mathrm{s}$):} The beverage carton rests upright on the tabletop at $(198.3, -134.5, -196.4)\,\mathrm{mm}$ while the right 7-DoF arm approaches from its home configuration ($d_{\mathrm{grip}}$ decreasing from $649.0\,\mathrm{mm}$ to $81.6\,\mathrm{mm}$).
    \item \textbf{Phase II --- Parallel-Jaw Grasp Closure (Frames 140--189, $10.60$--$14.27\,\mathrm{s}$):} The right parallel-jaw fingers enclose the carton body, reaching physical grasp contact at \textbf{Frame 162 ($t = 12.20\,\mathrm{s}$)} with a 3D pad-midpoint-to-carton-centroid distance of $\mathbf{11.05\,\mathrm{mm}}$, independently validating the bilateral $\mathrm{SE}(3)^3$ kinematic closure during dynamic manipulation.
    \item \textbf{Phase III --- Tabletop Lift and Lateral Transport (Frames 190--225, $14.33$--$16.87\,\mathrm{s}$):} The operator lifts the carton to $+82.2\,\mathrm{mm}$ clearance above the cup rim plane and translates it laterally toward the target cup ($\Delta r_{xy}$ decreasing from $199.9\,\mathrm{mm}$ to $43.5\,\mathrm{mm}$).
    \item \textbf{Phase IV --- Sustained Bullseye Pouring Window (Frames 226--352, $16.93$--$26.33\,\mathrm{s}$, $124$ frames / $9.40\,\mathrm{s}$):} The pouring spout enters the $R_{\mathrm{cup}} = 42.7\,\mathrm{mm}$ aperture cylinder at \textbf{Frame 226} ($\Delta r_{xy} = 27.9\,\mathrm{mm}$, margin $m_{\mathrm{rim}} = +14.5\,\mathrm{mm}$, $\Delta z = +68.7\,\mathrm{mm}$, $\theta = 25.6^\circ$), reaches \textbf{peak bullseye alignment at Frame 275 ($t = 20.67\,\mathrm{s}$)} with $\mathbf{\Delta r_{xy} = 9.8\,\mathrm{mm}}$ ($\mathbf{m_{\mathrm{rim}} = +32.8\,\mathrm{mm}}$ inside the cup rim, $\Delta z = +40.5\,\mathrm{mm}$, $\theta = 63.2^\circ$), reaches \textbf{maximum pouring tilt at Frame 315 ($t = 23.47\,\mathrm{s}$)} at $\mathbf{\theta = 79.4^\circ}$ while remaining safely inside the rim ($m_{\mathrm{rim}} = +9.1\,\mathrm{mm}$, $\Delta z = +33.3\,\mathrm{mm}$), and exits the rim cylinder at \textbf{Frame 352} ($\theta = 49.2^\circ$).
    \item \textbf{Phase V --- Untilt, Retraction, and Tabletop Placement (Frames 353--497, $26.40$--$37.47\,\mathrm{s}$):} The operator untilts the carton back to upright ($\theta = 2.8^\circ$), places it back onto the tabletop support plane, and retracts the right arm.
\end{enumerate}

\begin{table*}[t]
    \centering
    \caption{Phase-Level Quantitative Telemetry Summary and Key Physical Milestones Across the 498-Frame ($37.47\,\mathrm{s}$) Bimanual Pouring Sequence}
    \label{tab:phase_summary}
    \scriptsize
    \resizebox{\textwidth}{!}{%
    \begin{tabular}{@{}llcccccccc@{}}
        \toprule
        \textbf{Manipulation Phase /} & \textbf{Frame Span} & \textbf{Time Span} & \textbf{Spout Position in $O_{\mathrm{table}}$} & \textbf{Radial Offset} & \textbf{Rim Margin} & \textbf{Clearance} & \textbf{Pour Tilt} & \textbf{Gripper Dist.} \\
        \textbf{Key Physical Milestone} & \textbf{(Count)} & $t$ (s) & $(X,\, Y,\, Z)$ (mm) & $\Delta r_{xy}$ (mm) & $m_{\mathrm{rim}}$ (mm) & $\Delta z$ (mm) & $\theta$ (deg) & $d_{\mathrm{grip}}$ (mm) \\
        \midrule
        \multicolumn{9}{@{}l}{\emph{Part A: Aggregated Trajectory Statistics Across the Five Manipulation Phases}} \\
        \quad Phase I: Pre-Grasp Reach & F000--F139 ($140$) & $0.00$--$10.53$ & $(198.4, -134.8, -196.8)$ mean & $208.9 \pm 0.4$ & $-166.3$ & $+34.2$ & $0.1^\circ$--$0.5^\circ$ & $649.0 \to 81.6$ \\
        \quad Phase II: Grasp Closure & F140--F189 ($50$) & $10.60$--$14.27$ & $(198.5, -134.8, -197.1)$ mean & $208.6 \pm 0.3$ & $-166.2$ & $+34.0$ & $0.5^\circ$--$1.1^\circ$ & $\mathbf{11.05}$ min \\
        \quad Phase III: Lift \& Transport & F190--F225 ($36$) & $14.33$--$16.87$ & $(207.4 \to 364.6, -118.2, -236.9)$ & $199.9 \to 43.5$ & $-157.4 \to -1.1$ & $+33.0 \to +82.2$ & $1.2^\circ \to 23.5^\circ$ & $28.2$ mean \\
        \quad \textbf{Phase IV: Bullseye Pouring} & \textbf{F226--F352 ($\mathbf{127}$)} & $\mathbf{16.93}$--$\mathbf{26.33}$ & $\mathbf{(384.6, -118.9, -218.4)}$ mean & $\mathbf{24.6}$ mean & $\mathbf{+18.1}$ mean & $\mathbf{+55.2}$ mean & $\mathbf{25.6^\circ}$--$\mathbf{79.4^\circ}$ & $\mathbf{32.6}$ mean \\
        \quad Phase V: Untilt \& Retract & F353--F497 ($145$) & $26.40$--$37.47$ & $(350.4 \to 168.9, -152.1, -192.8)$ & $57.8 \to 239.6$ & $-15.3 \to -196.5$ & $+135.3 \to +31.0$ & $42.4^\circ \to 2.8^\circ$ & $35.8 \to 664.1$ \\
        \midrule
        \multicolumn{9}{@{}l}{\emph{Part B: Instantaneous State at Six Key Physical Manipulation Milestones}} \\
        \quad \textbf{M1: Grasp Closure Peak} & \textbf{Frame 162} & $\mathbf{12.20}$ & $\mathbf{(+198.6, -135.0, -197.4)}$ & $208.5$ & $-166.2$ & $+34.7$ & $0.5^\circ$ & $\mathbf{11.05}$ \\
        \quad \textbf{M2: Bullseye Rim Entry} & \textbf{Frame 226} & $\mathbf{16.93}$ & $\mathbf{(+379.4, -128.9, -232.0)}$ & $\mathbf{27.9}$ & $\mathbf{+14.5}$ & $\mathbf{+68.7}$ & $\mathbf{25.6^\circ}$ & $28.36$ \\
        \quad \textbf{M3: Sustained Pouring} & \textbf{Frame 245} & $\mathbf{18.33}$ & $\mathbf{(+381.7, -126.9, -244.4)}$ & $\mathbf{25.7}$ & $\mathbf{+17.0}$ & $\mathbf{+81.7}$ & $\mathbf{36.3^\circ}$ & $30.21$ \\
        \quad \textbf{M4: Peak Bullseye Align} & \textbf{Frame 275} & $\mathbf{20.67}$ & $\mathbf{(+397.6, -124.9, -203.6)}$ & $\mathbf{9.8}$ & $\mathbf{+32.8}$ & $\mathbf{+40.5}$ & $\mathbf{63.2^\circ}$ & $31.48$ \\
        \quad \textbf{M5: Maximum Pour Tilt} & \textbf{Frame 315} & $\mathbf{23.47}$ & $\mathbf{(+379.0, -109.4, -196.2)}$ & $\mathbf{33.4}$ & $\mathbf{+9.1}$ & $\mathbf{+33.3}$ & $\mathbf{79.4^\circ}$ & $34.13$ \\
        \quad \textbf{M6: Bullseye Rim Exit} & \textbf{Frame 352} & $\mathbf{26.33}$ & $\mathbf{(+369.6, -113.6, -260.0)}$ & $\mathbf{39.9}$ & $\mathbf{+2.8}$ & $\mathbf{+97.0}$ & $\mathbf{49.2^\circ}$ & $35.20$ \\
        \bottomrule
    \end{tabular}}
\end{table*}

\begin{table}[t]
    \centering
    \caption{Comparison of Operator Viewing Modalities for Bimanual Teleoperation Assistance (Frame 245 / 275 Pouring State)}
    \label{tab:modality_comparison}
    \scriptsize
    \resizebox{\columnwidth}{!}{%
    \begin{tabular}{@{}lccc@{}}
        \toprule
        \textbf{Spatial Assistance Metric} & \textbf{2D Camera} & \textbf{3D Persp.} & \textbf{Orthographic} \\
        & \textbf{Stream} & \textbf{Workspace} & \textbf{XY/XZ + HUD} \\
        \midrule
        Axial parallax bias $\delta_{\mathrm{parallax}}$ (F245) & $57.2\,\mathrm{mm}$ & $57.2\,\mathrm{mm}$ & $\mathbf{0.0\,\mathrm{mm}}$ \\
        Target cup rim visibility under arm occlusion & Occluded & Completed & \textbf{Completed} \\
        Direct planar radial margin $m_{\mathrm{rim}}$ readout & No & Foreshortened & \textbf{Exact ($+32.8\,\mathrm{mm}$)} \\
        Direct vertical clearance $\Delta z$ readout & No & Foreshortened & \textbf{Exact ($+40.5\,\mathrm{mm}$)} \\
        Viewpoint transition latency & $120\,\mathrm{ms}$ & Manual orbit & $\mathbf{0.0\,\mathrm{ms}}$ ($60\,\mathrm{fps}$) \\
        \bottomrule
    \end{tabular}}
\end{table}

\begin{table*}[t]
    \centering
    \caption{Comprehensive Empirical Benchmark of All 31 Perception, Depth, Pointmap, Segmentation, 6-DoF Tracking, Point-Tracking, Feature-Matching, and Calibration Frameworks Evaluated on Synchronized Humanoid Manipulation Recordings and Public Benchmarks (NVIDIA RTX 5090)}
    \label{tab:framework_audit}
    \scriptsize
    \resizebox{\textwidth}{!}{%
    \begin{tabular}{@{}llp{6.6cm}cc@{}}
        \toprule
        \textbf{Model Family} & \textbf{Evaluated Framework / Method} & \textbf{Empirical Performance on Humanoid Manipulation Recordings \& Benchmarks} & \textbf{Latency / VRAM} & \textbf{Architectural Verdict \& Role} \\
        \midrule
        Depth Sensor & \textbf{Calibrated Multi-Cam RGB-D} & $0.0\,\mathrm{mm}$ ref.; $251{,}182$ tabletop inliers; $\le 9.6\,\mathrm{mm}$ cross-cam residual & $3.5\,\mathrm{ms}$ / --- & \textbf{Adopted:} Primary metric depth anchor \\
        Depth Model & \textbf{Prior Depth Anything (Prior-DA)}~\cite{wang2025priorda} & $0.037\,\mathrm{m}$ chest $p_{95}$ ($0.11\,\mathrm{m}$ mean $p_{95}$); $4{,}102$ hole pixels filled & $1{,}005\,\mathrm{ms}$ / $5.0\,\mathrm{GB}$ & \textbf{Retained:} Offline depth-hole completion \\
        Depth Model & \textbf{Apple Depth Pro}~\cite{bochkovskii2025depthpro} & $0.83\,\mathrm{m}$ mean $p_{95}$ ($0.696\,\mathrm{m}$ chest, $1.419\,\mathrm{m}$ right wrist, $0.482\,\mathrm{m}$ head) & $305\,\mathrm{ms}$ / $8.5\,\mathrm{GB}$ & \textbf{Rejected:} Insufficient metric accuracy \\
        Depth Model & \textbf{Depth Anything V2 Small}~\cite{yang2024depthv2} & $1.61\,\mathrm{m}$ mean $p_{95}$ ($0.558\,\mathrm{m}$ chest, $2.516\,\mathrm{m}$ right wrist, $1.838\,\mathrm{m}$ head) & $25\,\mathrm{ms}$ / $0.42\,\mathrm{GB}$ & \textbf{Rejected:} Relative scale ambiguity \\
        Depth Model & \textbf{Depth Anything V2 Metric}~\cite{yang2024depthv2} & $79.4\,\mathrm{m}$ mean $p_{95}$ on close-range tabletop manipulation & $23\,\mathrm{ms}$ / $0.64\,\mathrm{GB}$ & \textbf{Rejected:} Uncalibrated metric head \\
        Depth Model & \textbf{UniDepthV2}~\cite{piccinelli2025unidepthv2} & $2.44\,\mathrm{m}$ mean $p_{95}$; $>80\%$ large-error rate on valid sensor overlap & $6.3\,\mathrm{ms}$ / $0.62\,\mathrm{GB}$ & \textbf{Rejected:} Severe wrist/chest drift \\
        3D Pointmap & \textbf{VGGT-1B Multiview}~\cite{wang2025vggt} & $0.47\,\mathrm{m}$ mean $p_{95}$ ($0.372\,\mathrm{m}$ chest, $0.510\,\mathrm{m}$ right wrist, $0.527\,\mathrm{m}$ head) & $33\,\mathrm{ms}$ / $16.5\,\mathrm{GB}$ & \textbf{Proposal Only:} Coarse multiview prior \\
        3D Pointmap & \textbf{SLAM3R} (Causal Window)~\cite{liu2025slam3r} & $0.40\,\mathrm{m}$ mean $p_{95}$ on causal chest video sequences & $28\,\mathrm{ms}$ / $11.1\,\mathrm{GB}$ & \textbf{Proposal Only:} $40\,\mathrm{cm}$ drift exceeds tolerance \\
        3D Pointmap & \textbf{DUSt3R} (Chest+Head Pairs)~\cite{wang2024dust3r} & $0.69\,\mathrm{m}$ mean $p_{95}$ on synchronized stereo camera pairs & $27\,\mathrm{ms}$ / $7.1\,\mathrm{GB}$ & \textbf{Superseded:} Inferior to calibrated RGB-D \\
        3D Pointmap & \textbf{MASt3R} (Metric Pairs)~\cite{leroy2024mast3r} & $1.79\,\mathrm{m}$ mean $p_{95}$ on cross-view tabletop camera pairs & $36\,\mathrm{ms}$ / $9.7\,\mathrm{GB}$ & \textbf{Rejected:} High cross-view metric error \\
        3D Pointmap & \textbf{CUT3R} (Streaming Pointmap)~\cite{wang2025cut3r} & Rotary-embedding CUDA kernel assertion on RTX 5090 & Blocked / --- & \textbf{Deferred:} Kernel incompatibility \\
        VLM Grounding & \textbf{Florence-2 Phrase Grounding}~\cite{xiao2024florence2} & $273.0\,\mathrm{px}$ median object center error on held-out manipulation frames & $110\,\mathrm{ms}$ / $3.2\,\mathrm{GB}$ & \textbf{Rejected:} Severe localization drift \\
        VLM Grounding & \textbf{Grounding DINO}~\cite{liu2023grounding} & $454/461$ depth-backed frames; $3.79\,\mathrm{px}$ $p_{75}$ with periodic refresh & $86\,\mathrm{ms}$ / $2.4\,\mathrm{GB}$ & \textbf{Adopted:} Zero-shot object discovery \\
        Video Seg. & \textbf{SAM~2.1 (Role-Isolated + GPU Overlap)}~\cite{ravi2024sam2} & $2{,}322$ calls; $+52$ cup masks; $5.72\,\mathrm{ms}$ CPU prep overlapped in $24.2\,\mathrm{ms}$ GPU wait & $100.2\,\mathrm{ms}$ / $4.8\,\mathrm{GB}$ & \textbf{Adopted:} Role-isolated mask tracker \\
        Hybrid Seg. & \textbf{YOLO11s-seg + SAM~2 + Flow}~\cite{jocher2024yolo11,farneback2003two} & $100\%$ coverage; $9.825\,\mathrm{px}$ center $p_{95}$; $0.9524$ IoU; $98.3\%$ at $30\,\mathrm{fps}$ & $32.1\,\mathrm{ms}$ (4-cam) & \textbf{Adopted:} High-rate 4-camera model \\
        Hybrid Seg. & \textbf{YOLO11n-seg Hybrid Variant}~\cite{jocher2024yolo11} & Exceeded $\le 10.0\,\mathrm{px}$ target center error $p_{95}$ accuracy threshold & $24.8\,\mathrm{ms}$ (4-cam) & \textbf{Rejected:} Insufficient mask precision \\
        6-DoF Track & \textbf{EvidenceScene (BT-Native Branch)}~\cite{wen2021bundletrack} & $539/539$ val; $34/38$ pouring cand; Laval Clock/Jar strict $191/136$; $-1.99\,\mathrm{GiB}$ RSS & $56.0$--$60.1\,\mathrm{ms}$ $p_{50}$ & \textbf{Adopted:} Primary 6-DoF pose owner \\
        6-DoF Track & \textbf{BundleSDF} (Neural SDF)~\cite{wen2023bundlesdf} & $768/768$ chest/head obs.; Bleach-hard $P_{10}/C_{10}=58.8\%/63.9\%$ at Frame 438 & $97\,\mathrm{ms}$ ($10.3\,\mathrm{fps}$) & \textbf{Retained:} Offline neural surface verifier \\
        6-DoF Track & \textbf{FoundationPose}~\cite{wen2024foundationpose} \& \textbf{Any6D}~\cite{lee2025any6d} & Single-anchor / few-view novel-object 6-DoF pose initialization & $65$--$120\,\mathrm{ms}$ & \textbf{Candidate:} Relocalization module \\
        6-DoF Track & \textbf{BIT} (Bidirectional Iterative)~\cite{song2025bit} & Joint geometry-pose refinement (SAM + SoftRas~\cite{liu2019softras}); single-RGB & Offline refiner & \textbf{Adapted:} Mutual refinement loss concept \\
        6-DoF Track & \textbf{6DOPE-GS} (Gaussian Splatting)~\cite{jin20256dopegs} & $3.5\,\mathrm{Hz}$ online tracking/reconstruction rate (below $15\,\mathrm{Hz}$ target) & $285\,\mathrm{ms}$ ($3.5\,\mathrm{Hz}$) & \textbf{Adapted:} Keyframe admission rules \\
        6-DoF Track & \textbf{KV-Tracker} (3D KV-Cache)~\cite{elich2026kvtracker} & $27$--$30\,\mathrm{fps}$ multi-view transformer KV cache with SAM~2.1 & $35\,\mathrm{ms}$ / $6.5\,\mathrm{GB}$ & \textbf{Benchmark:} Non-commercial license \\
        Point Track & \textbf{EvidenceScene (P2P Branch)}~\cite{point2pose2026,doersch2024bootstap} & YCB/Laval shadow + $F,G,R,M,S$ repairs ($M+S$ Jar AUC $0.750$ vs. BT $0.723$) & $14.7\,\mathrm{ms}$ / $8.5\,\mathrm{GB}$ & \textbf{Shadow Audited:} Conditional graph repairs \\
        Point Track & \textbf{TAPNext++}~\cite{zholus2025tapnext} \& \textbf{CoTracker3}~\cite{karaev2024cotracker3} & $1024$-frame horizon; TAPNext++ uses $2.53\,\mathrm{GB}$ weights; CoTracker3 NC & $23\,\mathrm{ms}$ / --- & \textbf{Superseded:} Replaced by BootsTAPIR \\
        Feature Match & \textbf{LoFTR} (Detector-Free)~\cite{sun2021loftr} & Adopted in BundleTrack ($34/38$ wrist cand, $195/148$ Clock/Jar strict); rejected for $K$ fit & $42\,\mathrm{ms}$ / $2.1\,\mathrm{GB}$ & \textbf{Adopted (BT):} Pairwise tracker matcher \\
        Feature Match & \textbf{EfficientLoFTR (FP32 / Full-MP)}~\cite{wang2024eloftr} & Clock/Jar strict $196$--$197$ / $146$--$147$ (Jar AUC $0.813$ vs. LoFTR $0.824$) & $92.8$ / $109.9\,\mathrm{ms}$ & \textbf{Shadow:} Faster $p_{95}$, slight Jar AUC drop \\
        Feature Match & \textbf{XFeat (Semi-Dense Matcher)}~\cite{potje2024xfeat} & Wrist pouring candidates drop $34 \to 31/38$; jump $p_{95}$ inflates $40.4 \to 73.0\,\mathrm{mm}$ & $42.1$ / $54.2\,\mathrm{ms}$ & \textbf{Rejected:} Sparse-depth candidate loss \\
        Feature Match & \textbf{ALIKED + LightGlue (Sparse)}~\cite{zhao2023aliked,lindenberger2023lightglue} & Collapses from $34$ to $2/38$ wrist candidates ($36$ retained-reference failures) & $45.6$ / $68.5\,\mathrm{ms}$ & \textbf{Rejected:} Keypoints miss $0.96\%$ valid depth \\
        Feature Match & \textbf{RoMa} (Robust Dense Matcher)~\cite{edstedt2024roma} & $0.270\%$ focal error, but $5.13\,\mathrm{px}$ principal point drift across views & $68\,\mathrm{ms}$ / $3.4\,\mathrm{GB}$ & \textbf{Rejected:} Unstable cross-cam intrinsics \\
        3D Geometry & \textbf{Open3D + $O(\text{runs})$ RLE + Fast RANSAC}~\cite{zhou2018open3d,besl1992icp,rusu2009fpfh,fischler1981ransac} & $251{,}182$ tabletop inliers; RLE box/centroid $10.38 \to 0.015\,\mathrm{ms}$; tag plane $4.47 \to 0.89\,\mathrm{ms}$ & $0.90\,\mathrm{ms}$ (CPU) & \textbf{Adopted:} Exact RLE \& plane fitting \\
        Calibration & \textbf{Bilateral $\mathrm{SE}(3)^3$ Closure} (Ours) & $1{,}648$ train / $413$ zero-refit held-out frames: $2.52\,\mathrm{mm}$ (L) \& $6.56\,\mathrm{mm}$ (R) $p_{50}$ & $1.40\,\mathrm{ms}$ FK & \textbf{Adopted:} 4-cam-to-kinematic closure \\
        \bottomrule
    \end{tabular}}
\end{table*}

\section{Comprehensive Benchmark of 31 Perception Frameworks and Ablation Studies}
\label{sec:benchmark}

\subsection{Systematic Evaluation of Depth Foundation Models, 3D Pointmaps, Segmenters, Trackers, and Matchers}
Before converging on our sensor-anchored multi-camera 3D reconstruction architecture, we conducted a systematic experimental campaign across \textbf{31 external frameworks and internal algorithmic variants} on an NVIDIA GeForce RTX~5090 workstation ($32\,\mathrm{GB}$ VRAM) using synchronized humanoid manipulation recordings and public benchmarks (YCBInEOAT, Laval Vicon motion capture, YCB-M, DexYCB). Rather than omitting negative or intermediate trials, Table~\ref{tab:framework_audit} documents the empirical metrics, latency/VRAM footprints, and architectural roles across all model families:
\begin{itemize}
    \item \textbf{Monocular and Sensor-Prior Depth Foundation Models:} Pure monocular depth estimators---including Depth Anything V2 Small~\cite{yang2024depthv2} ($1.61\,\mathrm{m}$ mean $p_{95}$), Depth Anything V2 Metric~\cite{yang2024depthv2} ($79.4\,\mathrm{m}$ mean $p_{95}$), UniDepthV2~\cite{piccinelli2025unidepthv2} ($2.44\,\mathrm{m}$ mean $p_{95}$, $>80\%$ large-error rate), and Apple Depth Pro~\cite{bochkovskii2025depthpro} ($0.83\,\mathrm{m}$ mean $p_{95}$, $305\,\mathrm{ms}$/frame)---cannot replace calibrated RGB-D depth for millimeter-level pouring alignment. Prior Depth Anything (Prior-DA)~\cite{wang2025priorda} conditioned on RealSense depth achieved $0.037\,\mathrm{m}$ $p_{95}$ on the chest camera (filling $4{,}102$ depth-hole pixels on Frame 31), but its $1{,}005\,\mathrm{ms}$/frame latency and $0.205\,\mathrm{m}$ wrist-camera $p_{95}$ restrict it to an offline depth-completion role.
    \item \textbf{Multi-View RGB-to-3D Pointmap Transformers:} Feed-forward pointmap models---VGGT-1B~\cite{wang2025vggt} ($0.47\,\mathrm{m}$ mean $p_{95}$, $16.5\,\mathrm{GB}$ VRAM), SLAM3R~\cite{liu2025slam3r} ($0.40\,\mathrm{m}$ mean $p_{95}$, $11.1\,\mathrm{GB}$ VRAM), DUSt3R~\cite{wang2024dust3r} ($0.69\,\mathrm{m}$ mean $p_{95}$), and MASt3R~\cite{leroy2024mast3r} ($1.79\,\mathrm{m}$ mean $p_{95}$)---capture coarse multi-view topology but exhibit $37$--$179\,\mathrm{cm}$ metric drift, while CUT3R~\cite{wang2025cut3r} failed with a custom rotary-embedding CUDA kernel assertion on RTX~5090.
    \item \textbf{Open-Vocabulary Grounding, Video Segmentation, and Distilled Student Runtime:} Phrase grounding with Florence-2~\cite{xiao2024florence2} suffered a $273\,\mathrm{px}$ median center error on held-out manipulation frames and was rejected. Grounding DINO~\cite{liu2023grounding} paired with role-isolated SAM~2.1~\cite{ravi2024sam2} achieved $0.987$--$0.991$ median IoU and $100\%$ pose coverage. For high-rate four-camera streaming, our distilled YOLO11s-seg~\cite{jocher2024yolo11} + staged SAM~2 + optical flow~\cite{lucas1981iterative,farneback2003two} hybrid runtime satisfied all nine accuracy and latency criteria ($9.825\,\mathrm{px}$ center $p_{95}$, $0.9524$ IoU, $32.14\,\mathrm{ms}$ 4-camera batch $p_{95}$), whereas the smaller YOLO11n-seg~\cite{jocher2024yolo11} exceeded the $10\,\mathrm{px}$ accuracy threshold.
    \item \textbf{6-DoF Object Pose Trackers and Long-Horizon Point Trackers:} BundleTrack~\cite{wen2021bundletrack} achieved $539/539$ ($100\%$) continuous tracking on our validation sequence ($56.0\,\mathrm{ms}$ $p_{50}$, $62.9\,\mathrm{ms}$ $p_{95}$), $498/498$ on our multi-camera composite pouring study, and $34/38$ camera-local candidates ($32/38$ registered surfaces) on the single-stream right-wrist pouring clip. BundleSDF~\cite{wen2023bundlesdf} was validated across $768/768$ dense chest/head observations ($9.63\,\mathrm{mm}$ depth residual, $10.3\,\mathrm{fps}$). We also evaluated FoundationPose~\cite{wen2024foundationpose}, Any6D~\cite{lee2025any6d}, BIT~\cite{song2025bit}, 6DOPE-GS~\cite{jin20256dopegs} ($3.5\,\mathrm{Hz}$), KV-Tracker~\cite{elich2026kvtracker} ($27$--$30\,\mathrm{fps}$), and Point2Pose~\cite{point2pose2026} paired with Online BootsTAPIR~\cite{doersch2024bootstap}.
    \item \textbf{Detector-Free, Semi-Dense, and Sparse Feature Matchers under Sparse Depth Support:} While LoFTR~\cite{sun2021loftr} and RoMa~\cite{edstedt2024roma} failed cross-camera intrinsic inference ($1.752\% / 10.52\,\mathrm{px}$ and $0.270\% / 5.13\,\mathrm{px}$ drift), LoFTR serves as the reliable pairwise correspondence engine inside BundleTrack. Shadow replacement trials inside our corrected native BundleTrack owner show that EfficientLoFTR~\cite{wang2024eloftr} (FP32 and Full-MP) lowers Laval tracking $p_{95}$ from $118.1$--$122.9\,\mathrm{ms}$ to $109.9$--$116.4\,\mathrm{ms}$ with strong Clock strict successes ($197$ vs. $195$) but slightly lower Jar ADD AUC ($0.813$ vs. $0.824$). On the $38$-frame right-wrist pouring clip---where in-mask valid depth median is only $0.958\%$---semi-dense XFeat~\cite{potje2024xfeat} cuts tracking $p_{50}/p_{95}$ from $60.14/83.36\,\mathrm{ms}$ to $42.05/54.24\,\mathrm{ms}$ but loses $3$ candidates ($34 \to 31$) and inflates jump $p_{95}$ from $40.40$ to $72.96\,\mathrm{mm}$, while sparse ALIKED+LightGlue~\cite{zhao2023aliked,lindenberger2023lightglue} collapses from $34$ to $2/38$ candidates because sparse 2D keypoints miss the $0.96\%$ valid depth pixels. All lossy matcher replacements are therefore rejected.
\end{itemize}

\begin{table}[t]
    \centering
    \caption{Counterfactual Architectural Ablation Results on Paired Manipulation Sequences}
    \label{tab:counterfactual_ablations}
    \scriptsize
    \begin{tabular}{@{}p{3.4cm}cc@{}}
        \toprule
        \textbf{Ablated Design Dimension} & \textbf{Counterfactual} & \textbf{Adopted Design} \\
        \midrule
        Depth back-projection mask margin & $+38.4\,\mathrm{mm}$ jump $p_{95}$ & $\mathbf{1\,\mathrm{px}}$ \textbf{erosion} ($0$ bleed) \\
        Object canonical gauge reference & Cumulative tilt drift & \textbf{Frame-0 gauge} ($100\%$) \\
        Multi-object SAM~2 memory bank & $-52$ masks, $+29$ FP & \textbf{Role-isolated} ($0$ leak) \\
        Aperture plane normal manifold & $11 \times 180^\circ$ flips & \textbf{Projective $\mathbb{RP}^2$} ($0$ flips) \\
        \bottomrule
    \end{tabular}
\end{table}

\subsection{Counterfactual Architectural Ablations}
Table~\ref{tab:counterfactual_ablations} quantifies four counterfactual design choices evaluated on paired replay sequences:
\begin{enumerate}
    \item \textbf{Foreground Mask Erosion ($1\,\mathrm{px}$) vs. Mask Dilation ($3$--$5\,\mathrm{px}$):} Dilating foreground segmentation masks by $3$--$5\,\mathrm{px}$ to increase point density causes depth pixels from the background tabletop or robot fingers to bleed across container boundaries, inflating inter-frame pose jump $p_{95}$ by $+38.4\,\mathrm{mm}$ and corrupting rim circle fits. Enforcing a $1\,\mathrm{px}$ foreground erosion margin eliminates boundary depth bleeding.
    \item \textbf{Frame-0 Canonical Object Gauge vs. Mid-Sequence Re-Anchoring:} Re-anchoring the canonical object frame $O_{\mathrm{obj}_0}$ mid-sequence using partial single-view depth support introduces cumulative rotational drift during high-tilt pouring ($\theta > 60^\circ$). Retaining the Frame-0 upright reference gauge with bounded mask-depth translation recovery (clipped to the $120.0\,\mathrm{mm}$ step bound) preserves $100\%$ pose continuity ($498/498$ frames).
    \item \textbf{Role-Isolated vs. Shared Multi-Object SAM~2 Memory:} Running primary and secondary object roles inside a shared SAM~2 predictor loses $52$ valid secondary cup masks and introduces $29$ background false-positive updates across $2{,}322$ calls compared to role-isolated sessions.
    \item \textbf{Projective $\mathbb{RP}^2$ vs. Directed $\mathbb{S}^2$ Rim Normals:} Treating raw SVD plane normals as directed vectors on $\mathbb{S}^2$ suffers $11$ discontinuous $180^\circ$ flips across $1{,}992$ frames, whereas our $\mathbb{RP}^2$ formulation reduces maximum temporal normal change to $3.11^\circ$ with $0.000\,\mathrm{mm}$ task-metric drift.
\end{enumerate}

\subsection{Factor-Graph Observability, Unanchored Gauge Drift, and the Availability--Surface Paradox ($F, G, R, M, S$)}
\label{sec:observability_audit}
To rigorously evaluate whether long-horizon 2D point trajectories (Point2Pose~\cite{point2pose2026} pinned to \texttt{51856226610d}) can replace our pairwise BundleTrack + \texttt{CausalObjectSurface} owner under partial occlusion and degenerate depth support, we conducted a paired mechanism study across six YCBInEOAT sequences and three Vicon-referenced Laval motion-capture sequences ($704$ frames, $71$ exact native GTSAM factor-graph replays, Table~\ref{tab:p2p_bt_observability}):
\begin{enumerate}
    \item \textbf{Metric-Dependent Winners and the Availability--Surface Paradox (Fork $F$):} On YCBInEOAT, BundleTrack achieves higher $10\,\mathrm{mm}$ ADD-S counts on all six sequences (e.g., Sugar $794$ vs. $617$, Bleach-hard $285$ vs. $257$), whereas Point2Pose achieves higher $0$--$100\,\mathrm{mm}$ recall AUC on Cracker ($0.822$ vs. $0.768$) and Bleach-hard ($0.720$ vs. $0.708$), and switching from nearest-vertex ADD-S to corresponding-vertex ADD on Sugar reverses the $10\,\mathrm{mm}$ winner ($206$ vs. $174$). On Laval Cookiejar, repairing two upstream Point2Pose defects in Fork $F$---historical keyframe pose overwrite ($f_{\mathrm{kf}} < f_t$) in \texttt{ModularPipeline.step} and pre-refinement pose return in \texttt{SVDClusterRANSACRegister.register}---raises pose availability from $207$ to $239/295$ and ADD AUC from $0.574$ to $0.635$, but \emph{reduces strict $20\,\mathrm{mm}$ origin successes from $136$ to $120$} and drops Frame-150 $10\,\mathrm{mm}$ surface model coverage $C_{10}$ from $76.82\%$ to $61.79\%$ ($P_{10} = 47.48\%$, vs. BundleTrack causal surfels at $P_{10}/C_{10} = 64.12\% / 99.71\%$).
    \item \textbf{Unanchored $\mathrm{SE}(3)$ Gauge Freedom ($G$) and Nuisance-Coupled Local Rank Deficiency ($R$):} Auditing all $704$ Laval graph records reveals that the current pose lacks a factor path to the sole initial prior $\mathbf{X}_0$ in $28/33$ Cookiejar and $6/16$ Clock-full updates ($0/22$ on Clock-hard), allowing Levenberg-Marquardt damping~\cite{zhang2016degeneracy,xicp} to converge with non-increasing cost while disconnected components $C' \cap C_0 = \emptyset$ drift freely and overwrite $6{,}519$ and $235$ frontend map points. Furthermore, eliminating uncertain landmark increments $\delta\mathbf{y}$ via orthogonal projection $\bar{\mathbf{A}}_x = (\mathbf{I} - \mathbf{A}_y \mathbf{A}_y^+) \mathbf{A}_x$ reveals that $8/34$ connected updates have local pose rank $r_i = \operatorname{rank}(\bar{\mathbf{A}}_x) - \operatorname{rank}(\bar{\mathbf{A}}_{x,-i}) < 6$. Prior-connected admission ($G$) raises Clock-full and Cookiejar ADD AUC to $0.739$ and $0.668$ and partially recovers Cookiejar $C_{10}$ to $69.94\%$, while marginal rank gating ($R$) raises Cookiejar $10\,\mathrm{mm}$ successes from $62$ to $97$ (AUC $0.718$) at the cost of Clock-hard recall ($0.545$).
    \item \textbf{Feature-Birth Observation Ledger ($M$) and Causal Surface-History Rebuild ($S$):} Inserting $394/25/83$ missing feature-birth \texttt{BearingRangeFactor}s ($M$) eliminates all disconnected updates ($0$ unanchored poses) but degrades standalone trajectories due to raw birth depth noise. Combining birth factors with atomic causal TSDF history rebuild~\cite{bundlefusion} at revised keyframe pose versions ($M+S$) exhibits a strong positive factorial interaction across all three actions, raising Cookiejar ADD AUC to $\mathbf{0.750}$ (exceeding BundleTrack's $0.723$), though BundleTrack retains higher strict $10\,\mathrm{mm}$ counts ($109$ vs. $81$) and superior surface coverage ($C_{10} = 99.94\%$ with composed native repairs).
\end{enumerate}

\begin{table}[t]
    \centering
    \caption{Comparative Observability and Surface Audit on Vicon-Referenced Laval Sequences (Clock-Hard $N=202$, Clock-Full $N=204$, Cookiejar $N=295$; Surface Snapshot at Frame 150 / Full)}
    \label{tab:p2p_bt_observability}
    \scriptsize
    \resizebox{\columnwidth}{!}{%
    \begin{tabular}{@{}lcccccc@{}}
        \toprule
        \textbf{Tracker / Mechanism Variant} & \textbf{Clock-H} & \textbf{Clock-F} & \textbf{Jar Avail.} & \textbf{Jar} & \textbf{Jar} & \textbf{Jar Surf.} \\
        & \textbf{AUC} & \textbf{AUC} & \textbf{/ $<10\text{mm}$} & \textbf{AUC} & \textbf{Unanch.} & $P_{10} / C_{10}$ (\%) \\
        \midrule
        Point2Pose Original ($P$)~\cite{point2pose2026} & $0.188$ & $0.716$ & $207 / 63$ & $0.574$ & $25/29$ & $55.6 / 76.8$ \\
        EvidenceScene P2P-$F$ (Correctness) & $0.658$ & $0.721$ & $239 / 62$ & $0.635$ & $28/33$ & $47.5 / 61.8$ \\
        EvidenceScene P2P-$G$ (Connected) & $0.658$ & $\mathbf{0.739}$ & $239 / 62$ & $0.668$ & $28/33$ & $50.2 / 69.9$ \\
        EvidenceScene P2P-$R$ (Rank Gate) & $0.545$ & $0.715$ & $241 / \mathbf{97}$ & $0.718$ & $28/34$ & $46.9 / 67.3$ \\
        EvidenceScene P2P-$M$ (Birth Ledger) & $0.207$ & $0.619$ & $231 / 53$ & $0.613$ & $\mathbf{0/30}$ & $55.4 / 67.1$ \\
        EvidenceScene P2P-$G+S$ & $0.453$ & $0.706$ & $229 / 66$ & $0.642$ & $25/30$ & $54.1 / 62.7$ \\
        \textbf{EvidenceScene P2P-$M+S$} & $0.614$ & $0.662$ & $231 / 81$ & $\mathbf{0.750}$ & $\mathbf{0/28}$ & $44.6 / 54.0$ \\
        \midrule
        BundleTrack Baseline~\cite{wen2021bundletrack} & $0.547$ & $0.880$ & $295 / 109$ & $0.723$ & N/A & $64.1 / 99.7$ \\
        \textbf{EvidenceScene BT-Native (Full)} & --- & $\mathbf{0.913}$ & $\mathbf{295 / 136^*}$ & $\mathbf{0.808}$ & N/A & $\mathbf{65.0 / 99.9}$ \\
        \bottomrule
    \end{tabular}}
    \vspace{1pt}
    {\tiny $^*$EvidenceScene BT-Native rows report strict ($\le 10\,\mathrm{mm}$ and $\le 10^\circ$) counts and full-sequence $4\,\mathrm{mm}$-grid surface metrics.}
\end{table}

\section{Discussion, Claim Boundaries, and Limitations}
\label{sec:discussion}
While our automated 3D reconstruction-driven teleoperation assistance system resolves core workspace completion, bilateral kinematic closure, and orthographic depth-perception challenges on real-robot manipulation data, we enforce strict provenance and deployment claim boundaries:
\begin{enumerate}
    \item \textbf{Provenance Separation Across Composite, Single-Stream, and Live Routes:} The 498-frame multi-camera pouring analysis (\texttt{20260804\_170929}) evaluates an offline-synthesized composite scene reconstruction ($498/498$ pose trajectory, $5{,}699$ accumulated surfels, and watertight parametric cup/carton completion). By contrast, on the single-stream configured pouring route (\texttt{20260625\_173830}, $38$ tracked frames after Frame-0 seed), the live-compatible \textbf{EvidenceScene} (BundleTrack + \texttt{CausalObjectSurface}) owner produces $34/38$ camera-local pose candidates and $32/38$ registered tabletop surfaces while explicitly reporting $4$ fail-closed \texttt{NO\_BA} frames rather than fabricating continuity.
    \item \textbf{Live Site-Acceptance Gate Status ($0/6$) and Remaining GPU Compute Gap:} The production contract requires live ROS-callback-to-trusted-depth-surface-render $p_{95} \le 100\,\mathrm{ms}$ across at least $300$ traces over $60\,\mathrm{s}$, one-hour accepted world updates $\ge 15\,\mathrm{Hz}$ overall and at interval $p_{05}$, readiness $\ge 99\%$, observation age $p_{95} \le 2\,\mathrm{s}$, one-hour RSS/CUDA growth $\le 64\,\mathrm{MiB}$, and physical verification of $2\,\mathrm{cm}$ grasp contact, $1\,\mathrm{cm}$ aperture/support height, $5\,\mathrm{mm}$ static drift, and $5\,\mathrm{px}$ HUD bounds. A read-only audit records $0/6$ live site-acceptance gates passed (\texttt{gate\_passed=false}, \texttt{wire\_authorized=false}) because the local ROS process loops an opened development bag rather than live camera hardware, and while our \texttt{2026-11-02} bit-exact adapters reduced warm HTTP \texttt{endpoint\_ms} from $176.40\,\mathrm{ms}$ ($155.33/219.06\,\mathrm{ms}$ $p_{50}/p_{95}$) to $\mathbf{127.80 / 153.40\,\mathrm{ms}}$ (cutting outer wrapper overhead by $72.9\%$ and eliminating $140.58\,\mathrm{ms}$ Generation-2 GC stalls), sequential SAM~2 + native BundleTrack GPU execution ($97.98 / 122.21\,\mathrm{ms}$) still requires TensorRT FP16 vision encoding to clear the $\le 100\,\mathrm{ms}$ end-to-end budget.
    \item \textbf{Decoupled Visual Assistance and Stationary Receptacle Scope:} Completed 3D meshes and orthographic HUD overlays inform the human operator without overriding low-level joint controllers, and the completed cup mesh assumes a stationary tabletop receptacle during pouring. Extending direction-selective joint pose-landmark optimization to bimanually moving receptacles under held-out physical evaluation is the primary research direction.
\end{enumerate}

\section{Conclusion}
\label{sec:conclusion}
We have presented \textbf{EvidenceScene} (\url{https://tiwater.github.io/evidencescene/}), an automated, competitor-agnostic 3D scene reconstruction-driven teleoperation assistance framework and open-source library (\texttt{evidencescene} v2.0.0) for humanoid bimanual manipulation. Grounded in a systematic empirical benchmark of \textbf{31 candidate foundation, matching, tracking, and reconstruction frameworks}, an exact factor-graph observability ($F,G,R,M$) and causal surface-transaction ($S$) audit across $704$ Vicon-referenced frames, and validation across $1{,}648$ training calibration frames, $413$ zero-refit held-out calibration frames, and $498$ four-camera manipulation frames ($1{,}992$ images) on a 61-link humanoid manipulator, \textbf{EvidenceScene} achieves $2.52\,\mathrm{mm}$ (left arm) and $6.56\,\mathrm{mm}$ (right arm) median bilateral kinematic closure, $1.60\,\mathrm{mm}$ dual-camera cup rim agreement, $191/204$ and $136/295$ strict public Laval tracking successes with $928$--$1{,}994\,\mathrm{MiB}$ RSS reduction under guarded frame retirement, $30.24 / 33.05\,\mathrm{ms}$ ($p_{50}/p_{95}$, $33.1\,\mathrm{fps}$) controlled YCBInEOAT pose-tracking throughput, an $8.65\times$ geometric-mean kernel speed-up ($93.0\times$ snapshot, $69.8\times$ update, $14.8\times$ budget selection) with $100\%$ bit-identical output digests, $85.1\%$ public and $90.0\%$ task $10\,\mathrm{mm}$ surface coverage at a $1{,}500$-point budget, $127.80 / 153.40\,\mathrm{ms}$ ($p_{50}/p_{95}$) bit-exact warm HTTP endpoint latency, $11.05\,\mathrm{mm}$ physical grasp-closure accuracy, and $9.8\,\mathrm{mm}$ peak pouring bullseye alignment ($+32.8\,\mathrm{mm}$ inside the cup rim at $63.2^\circ$ tilt) with $0\,\mathrm{ms}$ orthographic viewpoint switching in the operator console.

{\def\footnotesize{\fontsize{6.7pt}{7.5pt}\selectfont}
\setlength{\IEEEbibitemsep}{0pt}
\begin{thebibliography}{117}
\bibitem{chen2007human} J.~Y.~C. Chen, E.~C. Haas, and M.~J. Barnes, ``Human performance issues and user interface design for teleoperated robots,'' \emph{IEEE Trans. Syst., Man, Cybern. C}, vol.~37, no.~6, pp.~1231--1245, 2007.
\bibitem{rea2022still} D.~J. Rea and S.~H. Seo, ``Still not solved: A call for renewed focus on user-centered teleoperation interfaces,'' \emph{Frontiers in Robotics and AI}, vol.~9, p.~704056, 2022.
\bibitem{billard2019trends} A.~Billard and D.~Kragic, ``Trends and challenges in robot manipulation,'' \emph{Science}, vol.~364, no.~6446, p.~eaat8414, 2019.
\bibitem{chi2023diffusion} C.~Chi \emph{et al.}, ``Diffusion policy: Visuomotor policy learning via action diffusion,'' in \emph{Proc. Robotics: Science and Systems (RSS)}, 2023.
\bibitem{zhao2023learning} T.~Z. Zhao, V.~Kumar, S.~Levine, and C.~Finn, ``Learning fine-grained bimanual manipulation with low-cost hardware,'' in \emph{Proc. RSS}, 2023.
\bibitem{sheridan1992telerobotics} T.~B. Sheridan, \emph{Telerobotics, Automation, and Human Supervisory Control}. Cambridge, MA: MIT Press, 1992.
\bibitem{milgram1993argos} P.~Milgram, S.~Zhai, D.~Drascic, and J.~Grodski, ``Applications of augmented reality for human-robot communication,'' in \emph{Proc. IEEE/RSJ IROS}, 1993, pp.~1467--1472.
\bibitem{rakita2018autonomous} D.~Rakita, B.~Mutlu, and M.~Gleicher, ``An autonomous dynamic camera method for effective remote teleoperation,'' in \emph{Proc. ACM/IEEE HRI}, 2018, pp.~325--333.
\bibitem{rosenberg1993virtual} L.~B. Rosenberg, ``Virtual fixtures: Perceptual tools for telerobotic manipulation,'' in \emph{Proc. IEEE VRAIS}, 1993, pp.~76--82.
\bibitem{bowyer2013active} S.~A. Bowyer, B.~L. Davies, and F.~R. y~Baena, ``Active constraints/virtual fixtures: A survey,'' \emph{IEEE Trans. Robotics}, vol.~30, no.~1, pp.~138--157, 2014.
\bibitem{su2021mixed} Y.-H. Su \emph{et al.}, ``Mixed reality-enhanced multi-view teleoperation for robotic manipulation,'' \emph{IEEE Robot. Autom. Lett.}, vol.~7, no.~2, pp.~5461--5468, 2022.
\bibitem{schwarz2017nimbro} M.~Schwarz \emph{et al.}, ``NimbRo Picking: Versatile part handling for warehouse automation,'' in \emph{Proc. IEEE ICRA}, 2017, pp.~3032--3039.
\bibitem{fu2024mobile} Z.~Fu, T.~Z. Zhao, and C.~Finn, ``Mobile ALOHA: Learning bimanual mobile manipulation with low-cost whole-body teleoperation,'' in \emph{Proc. CoRL}, 2024.
\bibitem{qin2023anyteleop} Y.~Qin \emph{et al.}, ``AnyTeleop: A general vision-based dexterous robot arm-hand teleoperation system,'' in \emph{Proc. RSS}, 2023.
\bibitem{pi2026pi07} Physical Intelligence \emph{et al.}, ``$\pi_{0.7}$: A steerable generalist robotic foundation model with emergent capabilities,'' \emph{arXiv:2604.15483}, 2026.
\bibitem{yang2024depthv2} L.~Yang \emph{et al.}, ``Depth Anything V2,'' in \emph{Proc. NeurIPS}, 2024.
\bibitem{wang2025priorda} Z.~Wang \emph{et al.}, ``Prior Depth Anything: Marrying metric depth sensors with monocular depth foundation models,'' \emph{arXiv preprint}, 2025.
\bibitem{bochkovskii2025depthpro} A.~Bochkovskii \emph{et al.}, ``Depth Pro: Sharp monocular metric depth in less than a second,'' in \emph{Proc. ICLR}, 2025.
\bibitem{piccinelli2024unidepth} L.~Piccinelli \emph{et al.}, ``UniDepth: Universal monocular metric depth estimation,'' in \emph{Proc. IEEE/CVF CVPR}, 2024, pp.~10106--10116.
\bibitem{piccinelli2025unidepthv2} L.~Piccinelli \emph{et al.}, ``UniDepthV2: Universal monocular metric depth estimation made simpler,'' \emph{arXiv preprint}, 2025.
\bibitem{hu2024metric3dv2} M.~Hu \emph{et al.}, ``Metric3D v2: A versatile monocular geometric foundation model for zero-shot metric depth and surface normal estimation,'' \emph{IEEE Trans. Pattern Anal. Mach. Intell.}, 2024.
\bibitem{bhat2023zoedepth} S.~F. Bhat \emph{et al.}, ``ZoeDepth: Zero-shot transfer by combining relative and metric depth,'' \emph{arXiv:2302.12288}, 2023.
\bibitem{ke2024marigold} B.~Ke \emph{et al.}, ``Repurposing diffusion-based image generators for monocular depth estimation (Marigold),'' in \emph{Proc. IEEE/CVF CVPR}, 2024, pp.~9492--9502.
\bibitem{wang2024dust3r} S.~Wang, V.~Leroy, Y.~Cabon, B.~Chidlovskii, and J.~Revaud, ``DUSt3R: Geometric 3D vision made easy,'' in \emph{Proc. IEEE/CVF CVPR}, 2024, pp.~20697--20709.
\bibitem{leroy2024mast3r} V.~Leroy, Y.~Cabon, and J.~Revaud, ``Grounding image matching in 3D with MASt3R,'' in \emph{Proc. ECCV}, 2024.
\bibitem{wang2025vggt} J.~Wang \emph{et al.}, ``VGGT: Visual Geometry Grounded Transformer,'' in \emph{Proc. IEEE/CVF CVPR}, 2025.
\bibitem{liu2025slam3r} Y.~Liu \emph{et al.}, ``SLAM3R: Real-time dense scene reconstruction from monocular RGB videos,'' in \emph{Proc. IEEE/CVF CVPR}, 2025.
\bibitem{wang2025cut3r} Q.~Wang \emph{et al.}, ``Continuous 3D perception model with persistent state (CUT3R),'' in \emph{Proc. IEEE/CVF CVPR}, 2025.
\bibitem{wang2025spann3r} H.~Wang and L.~Agapito, ``3D reconstruction with spatial memory (Spann3R),'' in \emph{Proc. 3DV}, 2025.
\bibitem{zhang2025monst3r} J.~Zhang \emph{et al.}, ``MonST3R: A simple approach for estimating geometry in the presence of motion,'' in \emph{Proc. ICLR}, 2025.
\bibitem{liu2023grounding} S.~Liu \emph{et al.}, ``Grounding DINO: Marrying DINO with grounded pre-training for open-set object detection,'' in \emph{Proc. ECCV}, 2024.
\bibitem{xiao2024florence2} B.~Xiao \emph{et al.}, ``Florence-2: Advancing a unified representation for a variety of vision tasks,'' in \emph{Proc. IEEE/CVF CVPR}, 2024, pp.~4818--4829.
\bibitem{kirillov2023segment} A.~Kirillov \emph{et al.}, ``Segment anything,'' in \emph{Proc. IEEE/CVF ICCV}, 2023, pp.~4015--4026.
\bibitem{ravi2024sam2} N.~Ravi \emph{et al.}, ``SAM 2: Segment anything in images and videos,'' in \emph{Proc. ICLR}, 2025.
\bibitem{ren2024groundedsam} T.~Ren \emph{et al.}, ``Grounded SAM: Assembling open-world models for diverse visual tasks,'' \emph{arXiv:2401.14159}, 2024.
\bibitem{ren2024dinox} T.~Ren \emph{et al.}, ``DINO-X: A unified vision model for open-world object detection and understanding,'' \emph{arXiv:2411.14347}, 2024.
\bibitem{jocher2024yolo11} G.~Jocher, J.~Qiu, and A.~Chaurasia, ``Ultralytics YOLO11: Instance segmentation and real-time object detection,'' Ultralytics, 2024.
\bibitem{dosovitskiy2021vit} A.~Dosovitskiy \emph{et al.}, ``An image is worth 16x16 words: Transformers for image recognition at scale (ViT),'' in \emph{Proc. ICLR}, 2021.
\bibitem{ryali2023hiera} C.~Ryali \emph{et al.}, ``Hiera: A hierarchical vision transformer without the bells-and-whistles,'' in \emph{Proc. ICML}, 2023, pp.~29441--29454.
\bibitem{zhai2023siglip} X.~Zhai, B.~Mustafa, A.~Kolesnikov, and L.~Beyer, ``Sigmoid loss for language image pre-training (SigLIP),'' in \emph{Proc. IEEE/CVF ICCV}, 2023, pp.~11975--11986.
\bibitem{lucas1981iterative} B.~D. Lucas and T.~Kanade, ``An iterative image registration technique with an application to stereo vision,'' in \emph{Proc. IJCAI}, 1981, pp.~674--679.
\bibitem{farneback2003two} G.~Farneb\"ack, ``Two-frame motion estimation based on polynomial expansion,'' in \emph{Proc. SCIA}, 2003, pp.~363--370.
\bibitem{wen2021bundletrack} B.~Wen and K.~Bekris, ``BundleTrack: 6D pose tracking for novel objects without instance or category-level 3D models,'' in \emph{Proc. IEEE/RSJ IROS}, 2021, pp.~8067--8074.
\bibitem{wen2023bundlesdf} B.~Wen \emph{et al.}, ``BundleSDF: Neural 6-DoF tracking and 3D reconstruction of unknown objects,'' in \emph{Proc. IEEE/CVF CVPR}, 2023, pp.~606--617.
\bibitem{wen2024foundationpose} B.~Wen, W.~Yang, J.~Kautz, and S.~Birchfield, ``FoundationPose: Unified 6D pose estimation and tracking of novel objects,'' in \emph{Proc. IEEE/CVF CVPR}, 2024, pp.~17868--17879.
\bibitem{lee2025any6d} T.~Lee \emph{et al.}, ``Any6D: Model-free 6D pose estimation of novel objects,'' in \emph{Proc. IEEE/CVF CVPR}, 2025.
\bibitem{song2025bit} X.~Song \emph{et al.}, ``Prior-free 3D object tracking (BIT: Bidirectional Iterative Tracking),'' in \emph{Proc. IEEE/CVF CVPR}, 2025.
\bibitem{jin20256dopegs} Y.~Jin \emph{et al.}, ``6DOPE-GS: Online 6D object pose estimation using Gaussian Splatting,'' in \emph{Proc. IEEE/CVF ICCV}, 2025.
\bibitem{elich2026kvtracker} M.~Elich \emph{et al.}, ``KV-Tracker: Persistent 3D transformer key-value memory for real-time 6D object tracking,'' in \emph{Proc. IEEE/CVF CVPR}, 2026.
\bibitem{point2pose2026} ``Point2Pose: Causal multi-object 6-DoF tracking via long-term 2D point trajectories and incremental TSDF,'' \emph{arXiv:2604.10415}, 2026.
\bibitem{xiang2018posecnn} Y.~Xiang, T.~Schmidt, V.~Narayanan, and D.~Fox, ``PoseCNN: A convolutional neural network for 6D object pose estimation in cluttered scenes,'' in \emph{Proc. RSS}, 2018.
\bibitem{wang2019densefusion} C.~Wang \emph{et al.}, ``DenseFusion: 6D object pose estimation by iterative dense fusion,'' in \emph{Proc. IEEE/CVF CVPR}, 2019, pp.~3343--3352.
\bibitem{he2020pvn3d} Y.~He \emph{et al.}, ``PVN3D: A deep point-wise 3D keypoints voting network for 6DoF pose estimation,'' in \emph{Proc. IEEE/CVF CVPR}, 2020, pp.~11632--11641.
\bibitem{labbe2022megapose} Y.~Labb\'e \emph{et al.}, ``MegaPose: 6D pose estimation of novel objects via render \& compare,'' in \emph{Proc. CoRL}, 2022.
\bibitem{liu2022gen6d} Y.~Liu \emph{et al.}, ``Gen6D: Generalizable model-free 6-DoF object pose estimation from RGB images,'' in \emph{Proc. ECCV}, 2022, pp.~298--315.
\bibitem{he2022oneposepp} X.~He \emph{et al.}, ``OnePose++: Keypoint-free one-shot object pose estimation without CAD models,'' in \emph{Proc. NeurIPS}, 2022.
\bibitem{doersch2022tapvid} C.~Doersch \emph{et al.}, ``TAP-Vid: A benchmark for tracking any point in a video,'' in \emph{Proc. NeurIPS}, 2022.
\bibitem{doersch2023tapir} C.~Doersch \emph{et al.}, ``TAPIR: Tracking any point with per-frame initialization and temporal refinement,'' in \emph{Proc. IEEE/CVF ICCV}, 2023, pp.~10061--10072.
\bibitem{doersch2024bootstap} C.~Doersch \emph{et al.}, ``BootsTAP: Bootstrapped training for tracking-any-point (Online BootsTAPIR),'' in \emph{Proc. ECCV}, 2024.
\bibitem{zholus2025tapnext} A.~Zholus \emph{et al.}, ``TAPNext: Tracking any point over long video horizons,'' Google DeepMind, 2025.
\bibitem{karaev2024cotracker3} N.~Karaev \emph{et al.}, ``CoTracker3: Simpler and better point tracking by pseudo-labelling real videos,'' \emph{arXiv:2410.11831}, 2024.
\bibitem{lowe2004sift} D.~G. Lowe, ``Distinctive image features from scale-invariant keypoints (SIFT),'' \emph{Int. J. Comput. Vis.}, vol.~60, no.~2, pp.~91--110, 2004.
\bibitem{sarlin2020superglue} P.-E. Sarlin, D.~DeTone, T.~Malisiewicz, and A.~Rabinovich, ``SuperGlue: Learning feature matching with graph neural networks,'' in \emph{Proc. IEEE/CVF CVPR}, 2020, pp.~4938--4947.
\bibitem{lindenberger2023lightglue} P.~Lindenberger, P.-E. Sarlin, and M.~Pollefeys, ``LightGlue: Local feature matching at light speed,'' in \emph{Proc. IEEE/CVF ICCV}, 2023, pp.~17627--17638.
\bibitem{sun2021loftr} J.~Sun, Z.~Shen, Y.~Wang, H.~Bao, and X.~Zhou, ``LoFTR: Detector-free local feature matching with transformers,'' in \emph{Proc. IEEE/CVF CVPR}, 2021, pp.~8922--8931.
\bibitem{edstedt2024roma} J.~Edstedt \emph{et al.}, ``RoMa: Robust dense feature matching,'' in \emph{Proc. IEEE/CVF CVPR}, 2024, pp.~19790--19800.
\bibitem{curless1996volumetric} B.~Curless and M.~Levoy, ``A volumetric method for building complex models from range images,'' in \emph{Proc. ACM SIGGRAPH}, 1996, pp.~303--312.
\bibitem{newcombe2011kinectfusion} R.~A. Newcombe \emph{et al.}, ``KinectFusion: Real-time dense surface mapping and tracking,'' in \emph{Proc. IEEE ISMAR}, 2011, pp.~127--136.
\bibitem{izadi2011kinectfusion} S.~Izadi \emph{et al.}, ``KinectFusion: Real-time 3D reconstruction and interaction using a moving depth camera,'' in \emph{Proc. ACM UIST}, 2011, pp.~559--568.
\bibitem{niessner2013voxel} M.~Nie{\ss}ner, M.~Zollh\"ofer, S.~Izadi, and M.~Stamminger, ``Real-time 3D reconstruction at scale using voxel hashing,'' \emph{ACM Trans. Graph.}, vol.~32, no.~6, p.~169, 2013.
\bibitem{keller2013point} M.~Keller \emph{et al.}, ``Real-time 3D reconstruction in dynamic scenes using point-based fusion,'' in \emph{Proc. 3DV}, 2013, pp.~1--8.
\bibitem{whelan2015elasticfusion} T.~Whelan \emph{et al.}, ``ElasticFusion: Real-time dense SLAM and light source estimation,'' \emph{Int. J. Robot. Res.}, vol.~35, no.~14, pp.~1697--1716, 2016.
\bibitem{newcombe2015dynamicfusion} R.~A. Newcombe, D.~Fox, and S.~M. Seitz, ``DynamicFusion: Reconstruction and tracking of non-rigid scenes in real-time,'' in \emph{Proc. IEEE CVPR}, 2015, pp.~343--352.
\bibitem{dai2017bundlefusion} A.~Dai \emph{et al.}, ``BundleFusion: Real-time globally consistent 3D reconstruction using on-the-fly surface re-integration,'' \emph{ACM Trans. Graph.}, vol.~36, no.~3, p.~24, 2017.
\bibitem{schops2019badslam} T.~Sch\"ops, T.~Sattler, and M.~Pollefeys, ``BAD SLAM: Bundle adjusted direct RGB-D SLAM,'' in \emph{Proc. IEEE/CVF CVPR}, 2019, pp.~134--144.
\bibitem{zhou2018open3d} Q.-Y. Zhou, J.~Park, and V.~Koltun, ``Open3D: A modern library for 3D data processing,'' \emph{arXiv:1801.09847}, 2018.
\bibitem{park2019deepsdf} J.~J. Park, P.~Florence, J.~Straub, R.~Newcombe, and S.~Lovegrove, ``DeepSDF: Learning continuous signed distance functions for shape representation,'' in \emph{Proc. IEEE/CVF CVPR}, 2019, pp.~165--174.
\bibitem{mildenhall2020nerf} B.~Mildenhall \emph{et al.}, ``NeRF: Representing scenes as neural radiance fields for view synthesis,'' in \emph{Proc. ECCV}, 2020, pp.~405--421.
\bibitem{kerbl20233dgs} B.~Kerbl, G.~Kopanas, T.~Leimk\"uhler, and G.~Drettakis, ``3D Gaussian Splatting for real-time radiance field rendering,'' \emph{ACM Trans. Graph.}, vol.~42, no.~4, p.~139, 2023.
\bibitem{liu2019softras} S.~Liu, T.~Li, W.~Chen, and H.~Li, ``Soft Rasterizer: A differentiable renderer for image-based 3D reasoning,'' in \emph{Proc. IEEE/CVF ICCV}, 2019, pp.~7708--7717.
\bibitem{manuelli2019kpam} L.~Manuelli, W.~Gao, P.~Florence, and R.~Tedrake, ``kPAM: KeyPoint affordances for category-level robotic manipulation,'' in \emph{Proc. ISRR}, 2019, pp.~132--157.
\bibitem{qin2020keto} Z.~Qin, K.~Fang, Y.~Zhu, L.~Fei-Fei, and S.~Savarese, ``KETO: Learning keypoint representations for tool manipulation,'' in \emph{Proc. IEEE ICRA}, 2020, pp.~7278--7285.
\bibitem{florence2018dense} P.~R. Florence, L.~Manuelli, and R.~Tedrake, ``Dense object nets: Learning dense visual object descriptors by and for robotic manipulation,'' in \emph{Proc. CoRL}, 2018, pp.~373--385.
\bibitem{pan2022tax} C.~Pan, B.~Okorn, H.~Zhang, B.~Eisner, and D.~Held, ``TAX-Pose: Task-specific cross-pose estimation for robot manipulation,'' in \emph{Proc. CoRL}, 2022.
\bibitem{hartley2004mvg} R.~Hartley and A.~Zisserman, \emph{Multiple View Geometry in Computer Vision}, 2nd~ed. Cambridge Univ. Press, 2004.
\bibitem{besl1992icp} P.~J. Besl and N.~D. McKay, ``A method for registration of 3-D shapes (ICP),'' \emph{IEEE Trans. Pattern Anal. Mach. Intell.}, vol.~14, no.~2, pp.~239--256, 1992.
\bibitem{rusu2009fpfh} R.~B. Rusu, N.~Blodow, and M.~Beetz, ``Fast Point Feature Histograms (FPFH) for 3D registration,'' in \emph{Proc. IEEE ICRA}, 2009, pp.~3212--3217.
\bibitem{fischler1981ransac} M.~A. Fischler and R.~C. Bolles, ``Random sample consensus (RANSAC),'' \emph{Commun. ACM}, vol.~24, no.~6, pp.~381--395, 1981.
\bibitem{schonberger2016sfm} J.~L. Sch\"onberger and J.-M. Frahm, ``Structure-from-motion revisited (COLMAP),'' in \emph{Proc. IEEE CVPR}, 2016, pp.~4104--4113.
\bibitem{kaess2012isam2} M.~Kaess \emph{et al.}, ``iSAM2: Incremental smoothing and mapping using the Bayes tree,'' \emph{Int. J. Robot. Res.}, vol.~31, no.~2, pp.~216--235, 2012.
\bibitem{agarwal2022ceres} S.~Agarwal, K.~Mierle, and The Ceres Solver Team, ``Ceres Solver,'' Google Research, 2022.
\bibitem{hansen2001cmaes} N.~Hansen and A.~Ostermeier, ``Completely derandomized self-adaptation in evolution strategies (CMA-ES),'' \emph{Evol. Comput.}, vol.~9, no.~2, pp.~159--195, 2001.
\bibitem{pan2016robot} Z.~Pan, C.~Park, and D.~Manocha, ``Robot motion planning for pouring liquids,'' in \emph{Proc. ICAPS}, 2016, pp.~518--526.
\bibitem{schenck2017visual} C.~Schenck and D.~Fox, ``Visual closed-loop control for pouring liquids,'' in \emph{Proc. IEEE ICRA}, 2017, pp.~2629--2636.
\bibitem{kennedy2019autonomous} M.~Kennedy \emph{et al.}, ``Autonomous precision pouring from unknown containers,'' \emph{IEEE Robot. Autom. Lett.}, vol.~4, no.~3, pp.~2317--2324, 2019.
\bibitem{mardia2000directional} K.~V. Mardia and P.~E. Jupp, \emph{Directional Statistics}. Wiley, 2000.
\bibitem{taubin1991estimation} G.~Taubin, ``Estimation of planar curves, surfaces, and nonplanar space curves defined by implicit equations,'' \emph{IEEE Trans. Pattern Anal. Mach. Intell.}, vol.~13, no.~11, pp.~1115--1138, 1991.
\bibitem{tsai1989new} R.~Y. Tsai and R.~K. Lenz, ``A new technique for fully autonomous and efficient 3D robotics hand/eye calibration,'' \emph{IEEE Trans. Robot. Autom.}, vol.~5, no.~3, pp.~345--358, 1989.
\bibitem{park1994robot} F.~C. Park and B.~J. Martin, ``Robot sensor calibration: Solving AX=XB on the Euclidean group,'' \emph{IEEE Trans. Robot. Autom.}, vol.~10, no.~5, pp.~717--721, 1994.
\bibitem{horaud1995hand} R.~Horaud and F.~Dornaika, ``Hand-eye calibration,'' \emph{Int. J. Robot. Res.}, vol.~14, no.~3, pp.~195--210, 1995.
\bibitem{daniilidis1999hand} K.~Daniilidis, ``Hand-eye calibration using dual quaternions,'' \emph{Int. J. Robot. Res.}, vol.~18, no.~3, pp.~286--298, 1999.
\bibitem{andreff1999online} N.~Andreff, R.~Horaud, and B.~Espiau, ``On-line hand-eye calibration,'' in \emph{Proc. 3DIM}, 1999, pp.~430--436.
\bibitem{shah2019calibration} M.~Shah \emph{et al.}, ``Calibration of a multi-camera rig on a legged robot,'' in \emph{Proc. RSS}, 2019.
\bibitem{pretto2023eomp} A.~Pretto \emph{et al.}, ``EO@MP: A pose-graph approach to hand-eye calibration,'' in \emph{Proc. IEEE ICRA}, 2023.
\bibitem{pachtrachai2018handeye} K.~Pachtrachai \emph{et al.}, ``Hand-eye calibration without object: Simultaneous robot-world and hand-eye calibration,'' in \emph{Proc. IEEE/RSJ IROS}, 2018.
\bibitem{kalib2024} ``Kalib: Markerless hand-eye calibration with keypoint tracking,'' \emph{arXiv:2408.10562}, 2024.
\bibitem{jcr2024} ``JCR: Unifying scene representation and hand-eye calibration with 3D foundation models,'' \emph{arXiv:2404.11683}, 2024.
\bibitem{kcba2026} ``Robust camera-to-MoCap calibration via kinematic-chain bundle adjustment,'' \emph{arXiv:2604.22118}, 2026.
\bibitem{olson2011apriltag} E.~Olson, ``AprilTag: A robust and flexible visual fiducial system,'' in \emph{Proc. IEEE ICRA}, 2011, pp.~3400--3407.
\bibitem{wang2016apriltag2} J.~Wang and E.~Olson, ``AprilTag 2: Efficient and robust fiducial detection,'' in \emph{Proc. IEEE/RSJ IROS}, 2016, pp.~4193--4198.
\bibitem{krogius2019apriltag3} M.~Krogius, A.~Haggenmiller, and E.~Olson, ``Flexible layouts for fiducial tags (AprilTag 3),'' in \emph{Proc. IEEE/RSJ IROS}, 2019, pp.~1898--1903.
\bibitem{collins2014ippe} T.~Collins and A.~Bartoli, ``Infinitesimal plane-based pose estimation (IPPE),'' \emph{Int. J. Comput. Vis.}, vol.~109, no.~3, pp.~252--286, 2014.
\bibitem{sola2018micro} J.~Sol\`a, J.~Deray, and D.~Atchuthan, ``A micro Lie theory for state estimation in robotics,'' \emph{arXiv:1812.01537}, 2018.
\bibitem{barfoot2017state} T.~D. Barfoot, \emph{State Estimation for Robotics}. Cambridge Univ. Press, 2017.
\bibitem{huber1964robust} P.~J. Huber, ``Robust estimation of a location parameter,'' \emph{Ann. Math. Statist.}, vol.~35, no.~1, pp.~73--101, 1964.
\bibitem{macenski2022ros2} S.~Macenski \emph{et al.}, ``Robot Operating System 2: Design, architecture, and uses in the wild,'' \emph{Science Robotics}, vol.~7, no.~66, p.~eabm6074, 2022.
\bibitem{foxglove2022mcap} Foxglove Technologies, ``MCAP: Modular, high-performance serialization container format for heterogeneous pub/sub robotics data,'' 2022.
\bibitem{wang2024eloftr} Y.~Wang, X.~He, S.~Peng, D.~Tan, and X.~Zhou, ``EfficientLoFTR: Semi-dense local feature matching with sparse-like speed,'' in \emph{Proc. IEEE/CVF CVPR}, 2024, pp.~21666--21675.
\bibitem{potje2024xfeat} G.~Potje, F.~Cadar, A.~Araujo, R.~Martins, and E.~R. Nascimento, ``XFeat: Accelerated features for lightweight image matching,'' in \emph{Proc. IEEE/CVF CVPR}, 2024, pp.~2682--2691.
\bibitem{zhao2023aliked} X.~Zhao, X.~Wu, W.~Chen, P.~C.~Y. Chen, Q.~Xu, and Z.~Li, ``ALIKED: A lighter keypoint and descriptor extraction network via deformable transformation,'' \emph{IEEE Trans. Instrum. Meas.}, vol.~72, pp.~1--16, 2023.
\bibitem{zhang2016degeneracy} J.~Zhang, M.~Kaess, and S.~Singh, ``On degeneracy of optimization-based state estimation problems,'' in \emph{Proc. IEEE ICRA}, 2016, pp.~809--816.
\bibitem{xicp} T.~Tuna, J.~Nubert, Y.~Nava, S.~Khattak, and M.~Hutter, ``X-ICP: Localizability-aware LiDAR registration for robust localization in extreme environments,'' \emph{IEEE Trans. Robot.}, vol.~40, pp.~452--471, 2024.
\bibitem{bundlefusion} A.~Dai, M.~Nie{\ss}ner, M.~Zollh\"ofer, S.~Izadi, and C.~Theobalt, ``BundleFusion: Real-time globally consistent 3D reconstruction using on-the-fly surface re-integration,'' \emph{ACM Trans. Graph.}, vol.~36, no.~3, p.~24, 2017.
\end{thebibliography}}

\end{document}
